% concatenated from nov1.tex
\documentclass[a4paper]{article}

\usepackage{libertinus}
%\usepackage{datetime2}

\usepackage{xcolor}
\usepackage{graphicx}

\usepackage{amsmath,amsthm,amssymb,latexsym}

\newtheorem{theorem}{Theorem}[section]
\newtheorem{proposition}[theorem]{Proposition}
\newtheorem{lemma}[theorem]{Lemma}
\newtheorem{corollary}[theorem]{Corollary}
\newtheorem{remark}[theorem]{Remark}
\newtheorem{definition}[theorem]{Definition}
\newtheorem{example}[theorem]{Example}
\newtheorem{assumption}[theorem]{Assumption}
\newtheorem{conjecture}[theorem]{Conjecture}
\newtheorem{claim}[theorem]{Claim}

\newcommand{\indicator}[1]{\mathbf{1}_{\{ #1 \}}}


\definecolor{myblue}{RGB}{0, 102, 204}

\usepackage{bbm}
\usepackage{enumitem}
\usepackage{xcolor}

%\usepackage{showlabels}

\begin{document}

\title{On the utility problem in a market where price impact is transient}

\author{L\'or\'ant Nagy\thanks{HUN-REN Alfr\'ed R\'enyi Institute of Mathematics, Budapest, Hungary} \and Mikl\'os 
R\'asonyi\thanks{HUN-REN Alfr\'ed R\'enyi Institute of Mathematics and E\"otv\"os Lor\'and University, Budapest, 
Hungary; E-mail:\ rasonyi@renyi.hu}}

%\date{\protect\DTMnow}

\date{\today}

\maketitle

\begin{abstract}
We consider a discrete-time model of a financial market where a risky asset is bought and sold with transactions having
a transient price impact. It is shown that the corresponding utility maximization problem admits a solution. We manage to remove
some unnatural restrictions on the market depth and resilience processes that were present in earlier work. A non-standard 
feature of the problem is that the set of attainable portfolio values may fail the convexity property.
\end{abstract}

\noindent\textbf{Keywords:} price impact; market friction; optimal investment; dynamic programming; 
nonconvex domain of optimization

\noindent\textbf{MSC 2020:} Primary: 93E20, 91B70, 91B16; secondary: 91G10, 28B20

%\noindent\textbf{OR/MS subject classification:} Primary: utility/preference: theory, dynamic programming/optimal control; 
%secondary: finance: portfolio

\section{Introduction}

Investors' actions move market prices and make large position changes costly. More or less realistic models for
this \emph{price impact} phenomenon have been worked out in the mathematical finance literature.
One of the principal questions is the study of optimal investment in the presence of these (nonlinear)
transaction costs.

Price impact may be assumed \emph{instantaneous} if it affects the investor only at the moment of his/her portfolio rebalancing. 
This assumption leads to a relatively simple market dynamics,
see \cite{gr,pym}. At the other extreme, price impact may be \emph{permanent}, in which case
the investor's action pushes the price in a direction and this effect pertains to the whole future.
The most accurate description of reality is probably in between: price impact should
be \emph{transient}, with a gradually fading effect for the future, see \cite{peter-yan,bank-voss}.
The speed at which these effects disappear is characterized by \emph{market resilience}: if $r$ is
resilience then the bid-ask spread is diminished by a factor of $e^{-r}$ in one unit of time. The
size of the disturbance
caused by trading a unit amount of the asset is described by \emph{market depth}: if
$\delta$ is market depth then $1/\delta$ is the effect on the bid-ask spread of trading one unit of the risky asset. 
  
In the present article we prove that the uility maximization problem in discrete time for an agent
trading with transient price impact is well-posed: it admits a solution.
The problem under consideration has a complex, non-linear dynamics involving
all previous strategies at a given time. Moreover, the domain of optimization 
is \emph{non-convex}, which is a highly unusual feature. 

A continuous-time model with \emph{instantaneous} price impact was considered in \cite{gr}: 
they proved in their Theorem 5.1 that the utility maximization problem (with a concave utility function) admits
a solution under mild conditions. 

In \cite{bank-voss} a more advanced model with \emph{transient} 
price impact was treated where the markets' resiliance and depth were both assumed constant.
Theorem 3.3 of \cite{bank-voss} proved the existence of an optimal investment strategy in such a setting.

In \cite{peter-yan} market resilience and depth were both allowed to be stochastic processes
but a related monotonicity condition was imposed in their Assumption 2.4 which implies
that the set of attainable portfolio values is convex. That condition holds for
constant resilience and depth but otherwise it is rather restrictive. The paper \cite{peter-yan} did not
provide a general existence theorem for optimizers but a superreplication result (Theorem 3.2 in \cite{peter-yan})
with a dual characterization of contingent claims that can be superhedged from a given initial
position. They also proved a verification result (Corollary 3.5 in \cite{peter-yan}): a sufficient condition
implying that a given strategy is optimal.

In the present work we are dealing with the discrete-time version of the model of
\cite{peter-yan}. Our purpose is to prove the existence of an optimal strategy for
a utility maximizer while removing the stringent monotonicity assumption of \cite{peter-yan}
on market resilience and market depth, see Theorem \ref{main} and Remark \ref{megjegyzes} below. 

Section \ref{setup} presents the technical details of our model, the main result (Theorem \ref{main}) and some discussions
about the lack of convexity and its implications. Proofs will 
then be provided starting with Section \ref{poof}. Section \ref{single} deals with the one-step case while Section \ref{dyna} 
provides the dynamic programming argument. 
%Section \ref{aux} provides necessary technical lemmas. 
Further reflections are given in Section \ref{conclu}.

\section{Setup and results}\label{setup}

For $x\in\mathbb{R}$ we denote by $x^{+},x^{-}$ the positive and negative parts of $x$.
Fix a probability space $(\Omega,\mathcal{F},\mathbb{P})$ once and for all, together with a discrete-time
filtration $\mathcal{F}_{t}$, $t=0,\ldots,T$ where $\mathcal{F}_{0}$ coincides with $\mathbb{P}$-null sets. 
Mathematical expectation with respect to $\mathbb{P}$ will be denoted by $\mathbb{E}[\,\cdot\,]$,{}
$\mathbb{E}_{t}[\, \cdot\,]$ stands for $\mathcal{F}_{t}$-conditional expectation. $L^{0}$ is the set of all
real-valued random variables. $\mathbf{1}_{A}$ denotes the indicator function of a set $A$.

We now present the discrete-time version of the model in \cite{peter-yan}.
The time horizon of the investor will be some $T\in\mathbb{N}$. In the $T$th step the investor
must liquidate her position in the risky asset hence genuine decisions are made only at 
the previous times $t=1,\ldots, T-1$. To have a nontrivial problem we thus need to assume $T\geq 2$.

The risky asset's midprice (that is, the middle point of the bid-ask spread)
is described by an adapted real-valued process $P_{t}$, $t=0,\ldots,T$.{}
This is the price followed when there is no trading from the part of the investor in consideration.

Position in the risky asset at time $t$ is denoted by $X_{t}$, $t=0,\ldots,T$, we assume $X_{0}:=0$.
At each time $t$ the investor makes a portfolio adjustment based on previous information (before the new price $P_{t}$
is revealed) hence $X_{t}$ is assumed $\mathcal{F}_{t-1}$ measurable, that is, the strategy process
$X_{t}$, $1\leq t\leq T$ is predictable. We follow the convention $X_{-1}:=0$.
The set of all strategies is denoted by $\mathcal{A}$. We define
$$
\mathcal{A}_{0}:=\{X\in\mathcal{A}:X_{T}=0\},
$$
the set of strategies that liquidate the position in the risky asset by the end of the time horizon.
We note here, that due to the dynamics utilized in the paper, presented below, maximization of the utility of the terminal wealth $\xi_{T}^{X}$ in $X$ is economically meaningful only
over the set of strategies $\mathcal{A}_{0}$. Outside of $\mathcal{A}_0$, an investor could attain a position in the bank account with favourable expected utility while having large negative positions, and we would need to deal with liquidation value: such scenarios will be excluded.

%The actual midprice process $P^{X}_{t}$ is influenced by the agent's strategy so it follows the dynamics
%$$
%P^{X}_{t}=P_{t}+\sum_{j=1}^{t}\iota_{j} (X_{j}-X_{j-1}),\ t=0,\ldots,T
%$$
%where $\iota_{t}\geq 0$ is an adapted process describing the (permanent) impact of trading on the midprice.
%Note that \cite{peter-yan} had $\iota_{t}\equiv \iota$ constant, in which case simply $P_{t}^{X}=P_{t}+\iota X_{t}$.

We model liquidity with two primitives, resilience rate, and market depth: market resilience is described by a non-negative adapted process denoted by $r_{t}\geq 0$, $t=0,\ldots,T-1$, and market depth is a positive adapted process $\delta_{t}>0$, $t=1,\ldots,T$. The \emph{half-spread} follows a linear dynamics, namely
\begin{equation}\label{spred}
\zeta^{X}_{t+1}=e^{-r_{t}}\zeta^{X}_{t}+\frac{1}{\delta_{t+1}}|X_{t+1}-X_{t}|,\ 0\leq t\leq T-1,
\end{equation}
starting from an initial value $\zeta_{0}^{X}:=\zeta_{0}\geq 0$.
Finally, the cash account at time $t=1,\ldots,T$, considering that the investor pays the spread when trading, is calculated as
\begin{equation}\label{baccount}
\xi^{X}_{t+1}-\xi_{t}^{X}=-P_{t+1}(X_{t+1}-X_{t})-\zeta^{X}_{t+1}|X_{t+1}-X_{t}|,\ 0\leq t\leq T-1.
\end{equation}
Setting $\xi^{X}_{0}:=0$ for simplicity, with initial capital $z\in\mathbb{R}$, the investor has a wealth of $z+\xi_{T}^{X}$ at time $T$. The above model is the discrete-time counterpart of the model introduced in \cite{peter-yan}: except that in our setup only transient impact is modeled.


We further assume that the investor may possibly receive a random endowment during the trading period, 
described by an $\mathcal{F}_{T}$-measurable $\mathbb{R}$-valued
random variable $B$. Negative $B$ means that the investor has certain
payment obligations during the period considered.

The investor aims to maximize her expected utility from terminal wealth, hence we fix a \emph{utility function} $u:\mathbb{R}\to\mathbb{R}$.

\begin{assumption}\label{uint} The function $u$ is non-decreasing, continuous, $\lim_{x \to -\infty}u(x)=-\infty$, and $u$ is bounded from above.
We furthermore assume that for all $x,y,z\in\mathbb{R}$ and for all $t=1,\ldots,T-1$, 
\begin{equation}\label{inti}
\mathbb{E}[u(x+y P_{t}+zP_{T}-B)]>-\infty
\end{equation}
holds.
\end{assumption}

\begin{remark}{\rm If $u$ is concave then \eqref{inti} is implied by the simpler condition
\begin{equation*}
\mathbb{E}[u(x+y P_{t}-B)]>-\infty,\ t=1,\ldots,T.
\end{equation*}
Indeed, by concavity of the mapping $v\to u(x+v-B)$,
$$
\mathbb{E}\left[u(x+y P_{t}+zP_{T}-B)\right]\geq \frac{\mathbb{E}\left[u(x+2y P_{t}-B)\right]+\mathbb{E}\left[u(x+2zP_{T}-B)\right]}{2}.
$$
}
\end{remark}

Our next hypothesis stipulates that market depth is always above a fixed threshold.

\begin{assumption}\label{deltacska}
There is a constant $\delta>0$ such that $\delta_{t}\geq \delta$ almost surely, for all $t=1,\ldots,T$.	
\end{assumption}


Our main result asserts that an investor with an arbitrary initial capital $z\in\mathbb{R}$ may
find an optimal portfolio strategy $X^{*}(z)$.

\begin{theorem}\label{main} Let Assumption \ref{uint} and Assumption \ref{deltacska} be in force.
Then, for each $z\in\mathbb{R}$ there exists $X^{*}(z)\in \mathcal{A}_{0}$ such that
$$
\bar{u}(z):=\mathbb{E}\left[u\left(z+\xi^{X^{*}(z)}_{T}-B\right)\right]=\sup_{X\in\mathcal{A}_{0}}\mathbb{E}\left[u\left(z+\xi^{X}_{T}-
B\right)\right].
$$
\end{theorem}



Theorem \ref{main} shows that, despite the possible lack of convexity for the set of attainable values
(see Example \ref{notconvex} below), the utility maximization
problem admits an optimal strategy. We will present the proof of Theorem \ref{main} 
in the coming sections, using
a customized dynamic programming procedure.





From now on, for any stochastic process $Z_{t}$, we denote its increments by 
$\Delta Z_{t}:=Z_{t}-Z_{t-1}$, $0\leq t\leq T$ with the convention $Z_{-1}:=0$. 
It is convenient to use another parametrization for strategies: for a given real-valued
process $H_{t}$, $1\leq t\leq T$ such that $H_{t}$ is $\mathcal{F}_{t-1}$-measurable,
we may assign a unique strategy $X_{t}$ such that $X_{0}=0$, $\Delta X_{t}=H_{t}$, 
$1\leq t\leq T$. For such strategies we will also use the alternative notations 
$\zeta^{H},\xi^{H}$ for the corresponding half-spread and portfolio value processes.
With a slight abuse of notation we will also write $H\in \mathcal{A}_{0}$ when
we really mean that the corresponding $X$ is in $\mathcal{A}_{0}$. Note that $H\in\mathcal{A}_{0}$
implies that $H_{T}=-\sum_{j=1}^{T-1}H_{j}$, in particular, $H_{T}$ is $\mathcal{F}_{T-2}$-measurable.

Introduce the notation $\rho_{j,t}:=\exp\left[ -\sum_{i=j}^{t-1}r_{i}\right]$, $1\leq t\leq T$,
$0\leq j\leq t$. Note that $\rho_{t,t}=1$.
From \eqref{spred} and \eqref{baccount} we derive the explicit formula
\begin{equation}\label{expli}
\xi_{T}^{X}=-\sum_{t=1}^{T} P_{t}H_{t}
-\sum_{t=1}^{T} |H_{t}|\left( \rho_{0,t}\zeta_{0} +\sum_{j=1}^{t} 
\frac{\rho_{j,t}}{\delta_{j}}|H_{j}|\right).	
\end{equation}

\begin{example}\label{notconvex}{\rm{} Let $T=3$, $r_{t}=0$ for all $0\leq t\leq 3$, $P_{0}=\zeta_{0}=0$ and let 
$P_{1},P_{2},P_{3}$ be independent standard Gaussian.
Let $\mathcal{F}_{t}$, $0\leq t\leq 3$ be the natural filtration of the process $P$.
Let $\delta_{1}=1$, $\delta_{2}=\delta_{3}=10$. We claim that the set
$$
\mathcal{S}:=\{Y\in L^{0}:\exists X\in\mathcal{A}_{0}\mbox{ such that }\xi^{X}_{T}\geq Y\}
$$
fails convexity. We argue by contradiction. 
Convexity would imply, in particular, that for arbitrary \emph{non-negative deterministic} 
strategies $H_{i},G_{i}$,
$i=1,2$ such that $H_{3}:=-H_{1}-H_{2}$, $G_{3}=-G_{1}-G_{2}$, there would exist some $\bar{H}_{i}$, $i=1,2,3$ such that
$$
\xi^{\bar{H}}_{T}\geq \frac{\xi^{H}_{T}+\xi^{G}_{T}}{2}
$$
almost surely. In view of \eqref{expli}, this inequality can be rewritten as
\begin{eqnarray}\nonumber
& & \frac{1}{2}\left[H_{1}^{2}+H_{2}H_{1}+\frac{1}{10}H_{2}^{2}+(H_{1}+H_{2})H_{1}+\frac{1}{10}(H_{1}+H_{2})H_{2}+
\frac{1}{10}(H_{1}+H_{2})^{2}
\right]\\ 
\nonumber &+& \frac{1}{2}\left[G_{1}^{2}+G_{2}G_{1}+\frac{1}{10}G_{2}^{2}+(G_{1}+G_{2})G_{1}+\frac{1}{10}(G_{1}+G_{2})G_{2}
+\frac{1}{10}(G_{1}+G_{2})^{2}\right]\\
\nonumber &-& \left[\bar{H}_{1}^{2}+\bar{H}_{2}\bar{H}_{1}+\frac{1}{10}\bar{H}_{1}^{2}+(\bar{H}_{1}+\bar{H}_{2})
\bar{H}_{1}+\frac{1}{10}(\bar{H}_{1}+\bar{H}_{2})\bar{H}_{2}+
\frac{1}{10}(\bar{H}_{1}+\bar{H}_{2})^{2}\right] \\
\nonumber &\geq&	{}
\left(\bar{H}_{3}-\frac{H_{3}}{2}-\frac{G_{3}}{2}\right)P_{3}+\left(\bar{H}_{2}-\frac{H_{2}}{2}-\frac{G_{2}}{2}\right)P_{2}+
\left(\bar{H}_{1}-\frac{H_{1}}{2}-\frac{G_{1}}{2}\right)P_{1}.
\end{eqnarray}
Notice that the $\mathcal{F}_{2}$-conditional law of $(\bar{H}_{3}-H_{3}/2-G_{3}/2)P_{3}$
is nondegenerate Gaussian on the set $\{\bar{H}_{3}-H_{3}/2-G_{3}/2\neq 0\}$. Also, the left-hand
side and the last two terms of the right-hand side are $\mathcal{F}_{2}$ measurable. Hence
the above inequality necessarily implies that $\mathbb{P}(\bar{H}_{3}-H_{3}/2-G_{3}/2\neq 0)=0$.
By an analogous argument, also $\bar{H}_{i}-H_{i}/2-G_{i}/2=0$, $i=1,2$. But then 
\begin{eqnarray*}
& & \frac{1}{2}\left[\frac{21}{10}H_{1}^{2}+\frac{3}{10}H_{2}^{2}+\frac{23}{10}H_{1}H_{2}+
\frac{21}{10}G_{1}^{2}+\frac{3}{10}G_{2}^{2}+\frac{23}{10}G_{1}G_{2}\right]\\ 
&\geq& 
\frac{21}{10}\bar{H}_{1}^{2}+\frac{3}{10}\bar{H}_{2}^{2}+\frac{23}{10}\bar{H}_{1}\bar{H}_{2}.
\end{eqnarray*}
The latter property, however, badly fails. Take, for instance, $H_{1}=H_{2}=1$ and $G_{1}=1.5$, $G_{2}=0$.
We conclude that $\mathcal{S}$ is \emph{not} a convex set.}
%$$
%\frac{\phi(1,1)+\phi(1.5,0)}{2}<\phi(1.25,0.5),{}
%$$
%as easily checked. 
\end{example}

\begin{remark}\label{megjegyzes}{\rm The paper \cite{peter-yan} made a monotonicity
assumption. In the current discrete-time context it would require that the process $\rho_{0,t}^{2}\delta_{t}$ is a.s.
strictly decreasing in $t$. 
Under this hypothesis, \cite{peter-yan} showed that $\mathcal{S}$ is convex. The novel
contribution of our work is to drop such monotonicity assumptions and nevertheless to prove the existence 
of optimal strategies.}
\end{remark}




%\begin{eqnarray*}
%& & \frac{1}{2}\left[\frac{21}{10}H_{1}^{2}+\frac{3}{10}H_{2}^{2}+\frac{23}{10}H_{1}H_{2}+
%\frac{21}{10}H_{1}^{2}+\frac{3}{10}H_{2}^{2}+\frac{23}{10}H_{1}H_{2}\right]\\ 
%&\geq& 
%\frac{21}{10}\bar{H}_{1}^{2}+\frac{3}{10}\bar{H}_{2}^{2}+\frac{23}{10}\bar{H}_{1}\bar{H}_{2}.
%\end{eqnarray*}



We finally point out in a simple example why
non-convexity of the domain of optimization may cause trouble in optimal investment problems.

\begin{example}\label{noncon}
{\rm We consider a one-step frictionless market model. 
Let the family of permitted strategies be $\Phi:=\{0,1\}$. That is, the investor
may take either a unit position in the risky asset or no position at all. Let us consider the utility function 
$u(x):=\ln(x)$ for $x>0$. 
Let the investor have initial capital $z>1$. Assume that the return on his investment is
a random variable $R$ with  $\mathbb{P}(R=2)=1/2=\mathbb{P}(R=-1)$. His indirect utility is then 
$$
\bar{u}(z):=\sup_{\phi\in\Phi} \mathbb{E}[u(z+\phi R)]=\max\left\{\ln(z),{}
\frac{\ln(z+2)+\ln(z-1)}{2}\right\},\ z>1.
$$
This function fails concavity: it is non-differentiable
at $2$ with the right-hand derivative being strictly larger than the left-hand derivative.

We conclude that even if the investor's utility is risk-averse (concave), his/her \emph{indirect} utility
may well fail this property when the family of permitted strategies is non-convex. Hence
in related multistep optimization problems
one needs to deal with a dynamic programming procedure involving \emph{non-concave} 
functions, as in \cite{laurence}.}
\end{example}




\section{Preparation for the proof}\label{poof}

%Let $\Xi_{t}$ denote the set of $\mathcal{F}_{t}$-measurable real-valued random variables, $0\leq t\leq T$. 
%We will follow a more involved version of the dynamic programming approach 
%from \cite{bounded}.


%Define $M_{T}(n):=u(n-B)$, $n\in\mathbb{Z}$, these are integrable random variables by Assumption \ref{uint}.

Recall from \eqref{baccount} and \eqref{spred} that
$$
\Delta \xi_{t}^{H}=-P_{t} H_{t}-|H_{t}|
\left(\rho_{0,t}\zeta_{0}
+\sum_{j=1}^{t}  
\frac{\rho_{j,t}}{\delta_{j}}|H_{j}|\right).
$$
Inspired by this formula, for $1\leq t\leq T$ and $h = (h_{1},\ldots,h_{t})$ we introduce the random functions
\begin{equation}\label{kappafg}
\kappa_{t}(h):= 
-P_{t} h_{t}-|h_{t}|\left( \rho_{0,t}\zeta_{0}+\sum_{j=1}^{t} 
\frac{\rho_{j,t}}{\delta_{j}} |h_{j}|\right)
\end{equation}
for all $h_{1},\ldots,h_{t}\in\mathbb{R}$. Note that the mapping, describing the innovation corresponding to a deterministic action of the trader,
$$
h_{t}\to \kappa_{t}((h_{1},\ldots,h_{t-1},h_{t})),
$$
is \emph{concave} for every fixed $(h_{1},\ldots,h_{t-1})$, 
but $\kappa_{t}$, as a function of $t$ variables, has no reason to be concave.
Note also that innovation has an "action-independent" market bound, namely the quantity
\begin{equation}\label{lambda}
\kappa_{t}((h_{1},\ldots,h_{t}))\leq \lambda_{t}(h_{t}):=-P_{t} h_{t}
-\frac{h_{t}^{2}}{\delta_{t}} \leq \frac{P_t^2 \delta_t}{4}.
\end{equation}
%Finally we define $V_{T}:\Omega\times \mathbb{R}\times \mathbb{R}^{T}\to \mathbb{R}$ as
%\begin{align*}
%%\begin{equation}\label{lambda}
%%\kappa_{t}(h_{1},\ldots,h_{t})\leq \lambda_{t}(h_{1},\ldots,h_{t}):=-P_{t} h_{t}
%%-\frac{h_{t}^{2}}{\delta_{t}}
%%\end{equation}
%%Define $V_{T}:\Omega\times \mathbb{R}\times \mathbb{R}^{T}\to \mathbb{R}$ as
%%\begin{align*}
%V_{T}&(\omega,x, h_{1},\ldots,h_{T})
%\\
%:&=
%G_T(\omega, x + \kappa_T( h_1,\hdots, h_{T}), h_1,\hdots, h_{T}),\quad \omega\in\Omega,x, h_{1},\ldots, h_{T}\in\mathbb{R}.
%\end{align*}

We recall Lemma 6.8 of \cite{laurence}.

\begin{lemma}
\label{gast}
Let $(\Omega, {\cal H},P)$ be a complete probability space. Let $\Xi^d$ be the set of 
$\mathcal{H}$-measurable $d$-dimensional random variables.
Let $F:\Omega\times\mathbb{R}^d\to\mathbb{R}$ be a function such that for almost all $\omega\in\Omega$,
$F(\omega,\cdot)$ is continuous and for each $y\in\mathbb{R}^d$, $F(\cdot,y)$ is ${\cal H}$-measurable. Let $K>0$ be
an ${\cal H} $-measurable random variable.

Set $f(\omega)=\mathrm{ess.}\sup_{\xi\in \Xi^d, |\xi|\leq  K} F(\omega,\xi(\omega))$.  Then,
for almost all $\omega$,
\begin{eqnarray}
\label{lavieestbienfaite2}
f(\omega) & = & \sup_{y \in \mathbb{R}^d, |y| \leq K(\omega) } F(\omega,y).
\end{eqnarray}
\hfill $\square$
\end{lemma}

%\begin{proof}
%$F$ is easily seen to be $\mathcal{H}\otimes\mathcal{B}(\mathbb{R}^d)$-measurable and so is
%$$
%\sup_{y \in \mathbb{R}^d, |y| \leq K(\omega) } F(\omega,y)=\sup_{y \in \mathbb{Q}^d, |y| \leq K(\omega) } F(\omega,y).
%$$
%Hence $\sup_{y \in \mathbb{R}^d, |y| \leq K(\omega) } F(\omega,y) \geq  f(\omega)$ a.s.
%by the definition.
%Assume that the inequality is strict with positive probability. Then for some $\varepsilon>0$
%the set
%$$A=\{(\omega,y)\in \Omega \times  \mathbb{R}^d \ : \ |y| \leq K(\omega); \ F(\omega,y) -f(\omega)
%\geq \varepsilon \}$$
%has a projection $A'$ on $\Omega$ with $P(A')>0$.
%Recall that $\omega \rightarrow F(\omega,\xi(\omega))$ is ${\cal H}$-measurable for $\xi\in \Xi^d$.
%By definition of the essential supremum, $f$ is ${\cal H} $-measurable and
%hence $A \in {\cal H} \otimes {\cal B}(\mathbb{R}^d)$. The measurable selection theorem (see 
%Proposition III.44 in \cite{dm}) applies and
%there exists $\eta\in\Xi^d$ such that $(\omega, \eta(\omega))\in A$
%for $\omega\in A'$ (and  $\eta(\omega)=0$ on the complement of $A'$). This leads to a contradiction since for all $\omega \in A'$,
%$f(\omega)<F(\omega,\eta(\omega))$ by the construction of $\eta$ and
%$f(\omega)\geq
%F(\omega,\eta(\omega))$ a.s. by the definition of $f$.
%\end{proof}

A compactness result involving random subsequences comes next, this is Lemma 2 of \cite{ks}.

\begin{lemma}\label{randomsub} Let $\mathcal{H}\subset\mathcal{F}$ be a sigma-algebra. 
Let $X_n$ be a sequence of $\mathcal{H}$-measurable $d$-dimensional random variables such that
\begin{equation}\label{raffaello}
\liminf_{n\to\infty}|X_n|<\infty
\end{equation}
almost surely. Then there exist
$\mathcal{H}$-measurable random variables 
$n_k:\Omega\to\mathbb{N}$, $k\in\mathbb{N}$ with $n_k(\omega)<n_{k+1}(\omega)$ for all $\omega\in\Omega$
and $k\in\mathbb{N}$ and an $\mathcal{H}$-measurable 
random variable $X$ such that $X_{n_k}\to X$ a.s. In such a case we write that 
\emph{there exists an $\mathcal{H}$-measurable
random subsequence}. \hfill $\square$
\end{lemma}


The following lemma uses the same idea as Lemma A.3 of \cite{lukasz}; it provides continuous versions
for certain random fields.

\begin{lemma}\label{szeder} Let $\mathcal{H}\subset\mathcal{F}$ be a sigma-algebra.
Define $\mathcal{K}:=[-N,N]^{n}$.
Let $L:\Omega\times \mathcal{K}\to \mathbb{R}$ be such that for a.e. $\omega\in\Omega$,
$L(\omega,\cdot)$ is continuous and for all $x\in \mathcal{K}$, $L(\cdot,x)$ is
measurable such that $\sup_{z\in \mathcal{K}} |L(\omega,z)|$ is integrable.
Then there is $l:\Omega\times \mathcal{K}\to \mathbb{R}$ such that for a.e. $\omega\in\Omega$,
$l(\omega,\cdot)$ is continuous and for all $k\in \mathcal{K}$, $E(L(k)|\mathcal{H})=l(k)$ a.s.
\end{lemma}
\begin{proof}
Consider the separable Banach space 
$\mathbb{B}:=C([-N,N]^{n})$ of continuous functions on $[-N,N]^{n}$ with the supremum norm. Clearly, $L:\Omega\to \mathbb{B}$ and
for all $\mu$ in the dual space $\mathbb{B}'$ (which can be represented as a Borel signed measure), 
$\mu(L)=\int_{\mathcal{K}} L(\omega,x)\mu(dx)$
is a measurable function on $(\Omega,\mathcal{F})$: indeed, this is clear for
$\mu$ with finite support and then follows for general $\mu$ by approximation.
Note that, for each $k\in \mathcal{K}$, the linear functional
$f_k(x):=x(k)$, $x\in \mathbb{B}$ is continuous (w.r.t. the norm of $\mathbb{B}$) so
$f_k\in\mathbb{B}'$. Now it follows
from Proposition V.2.5. of \cite{neveu} that there is a measurable $l:\Omega\to \mathbb{B}$
such that $$
l(k)=f_k(l)=E(f_k(L)|\mathcal{H})=E(L(k)|\mathcal{H}),
$$ 
for each $k\in \mathcal{K}$, as claimed.
\end{proof}





Now we turn to a set of Lemmas that ensure that we can perform a backward iteration, 
and produce a series of actions that forms a candidate strategy of optimal execution.



\section{Single step case}\label{single}

Let $t\geq 1$ be an integer and let $\mathcal{G}, \mathcal{H}$ be $\mathbb{P}$-complete sigma-algebras 
over $\Omega$ such that $\mathcal{H} \subset \mathcal{G} \subset \mathcal{F}$ holds, and denote the set of $\mathcal{H}$-measurable 
$\mathbb{R}$-valued random variables by $\Xi$. We will consider functions
$$G_0:\Omega\times \mathbb{R} \times \mathbb{R}^{t}  \to \mathbb{R}, \ (x,v) \mapsto G_0(x,v),$$
and
$$V:\Omega\times \mathbb{R} \times \mathbb{R}^{t-1} \times \mathbb{R} \to \mathbb{R}, \ (x,j,h) \mapsto G_0(x + \kappa_{t}((j,h)),(j,h)),$$
where $(j,h) = (j_1,\hdots,j_{t-1},h)$, and the $\kappa_t$ is as in (\ref{kappafg}). Below we introduce assumptions that serve as a basis 
for the iterative arguments later.

\begin{assumption}\label{A1}
The function $G_0$ is 
$\mathcal{G}\otimes\mathcal{B}(\mathbb{R})\otimes\mathcal{B}(\mathbb{R}^{t})$-measurable, the mapping $G_0(\omega,\cdot, \cdot)$, is jointly continuous 
and non-decreasing in its first variable, $\mathbb{P}$-almost surely. 
%Furthermore for every $(x_0,v_0) \in \mathbb{R} 
%\times \mathbb{R}^{t}$ there exists an $\mathcal{G}$-measurable 
%$(\varepsilon_1,\hdots,\varepsilon_t)=\varepsilon = \varepsilon(\omega)\in\mathbb{R}^{t-1}$ so that on the region 
%$$\mathbb{R}\times [v_1-\varepsilon_1(\omega), v_1+\varepsilon_1(\omega)]\times \hdots \times [v_t-\varepsilon_t(\omega), 
%v_t+\varepsilon_t(\omega)]$$ the function $G_0(\omega,\cdot,\cdot)$ is either non-decreasing or non-increasing.
\end{assumption}

\begin{assumption}\label{A2}
There exists a function $\bar G_0:\Omega\times\mathbb{R} \mapsto \mathbb{R}$, and a constant $C >0$, such that for all $x \in \mathbb{R}$ there exists a zero measure set outside of which 
$$\bar G_0(x) \to - \infty$$
holds as $x \to -\infty$. Furthermore, $G_0(\omega,\cdot)$ is non-decreasing almost surely, and for all pairs $(x, v) \in \mathbb{R} \times \mathbb{R}^{t}$ we have that the inequalities
\begin{equation*}
G_0(x,v) \leq \bar G_0(x) \leq C
\end{equation*}
hold, again outside a set of measure zero.
%Furthermore, $G_0(x,v) \leq C$ for all $x \in \mathbb{R}$, $v \in \mathbb{R}^{t}$ and for almost every $\omega\in\Omega$.
\end{assumption}

%\begin{assumption}\label{A2}
%Assume, that there exists $\mathcal{H}\times\mathcal{B}(\mathbb{R})$-measurable $J:\Omega\times\mathbb{R}\to\mathbb{R}$ so that for all $x \in \mathbb{R}$ we have 
%\begin{equation*}
%V(x,h_1,\hdots,h_{t}) \leq  J(x)
%\end{equation*}
%almost surely for all $h_1, \hdots, h_{t} \in \mathbb{R}$.
%\end{assumption}

\begin{assumption}\label{A3}
Assume that for any $m \in \mathbb{N}$ and for all $1\leq t^{'} \leq t$, there exists an integrable random variable $M = M(m, t^{'})$ so that
$$M \leq G_0(x + \kappa_{t^{'}}((v_1,\hdots, v_{t^{'}})),v)$$
holds for every $x \in [-m,m]$, $v \in [-m,m]^{t}$, and for almost every $\omega \in \Omega$. 
%Furthermore, assume that there exists an integrable random variable $\bar M = M(m, t^{'})$ so that
%$$\bar M \leq \bar G_0(x + P_{t^{'}}^2 \delta_{t^{'}} / 4)$$
%holds for every $x  \in [-m,m]$ and for almost every $\omega \in \Omega$.
\end{assumption}



%\begin{remark}
%When Assumption \ref{A2} and Assumption \ref{A3} is in place, the lower bound $M$ -- in Assumption \ref{A3} -- for $G_0$ is a bound for  $\bar G_0$ also, with setting $v = 0 \in \mathbb{R}^{t}$ and $t^{'}=t$.
%\end{remark}

%\begin{assumption}\label{A3}
%Assume that for any compact set $M \in \mathbb{R}^{t+1}$ -- denoting with $\partial M$ its boundary -- there exists an integrable random variable $\hat M = \hat M(\partial M)$ so that $\hat M \leq V(x,j,h)$ and $\hat M \leq \bar G_0(x)$ holds for every $(x,j,h) \in M$ and for almost every $\omega \in \Omega$.
%\end{assumption}


Throughout Section \ref{single} we postulate that the conditions prescribed by Assumption \ref{A1}, Assumption \ref{A2}, and Assumption \ref{A3} hold.

\begin{lemma}\label{continuous_version}
There exists $L : \Omega\times \mathbb{R} \times \mathbb{R}^{t-1} \times \mathbb{R} \to \mathbb{R}$ so that for all $(x,j,h) \in \mathbb{R} \times \mathbb{R}^{t-1} \times \mathbb{R}$ we have 
$$ E[V(x,j,h) | \mathcal{H}] = L(x,j,h)$$
almost surely, and furthermore, $L(\cdot,\cdot,\cdot)$  is continuous in all its variables for almost all $\omega \in \Omega$.
\end{lemma}
 
\begin{proof}
Under Assumption \ref{A1}, Assumption \ref{A2}, and Assumption \ref{A3} the proof is a direct consequence of Lemma \ref{szeder}.

\end{proof}

\begin{lemma}\label{lemma::infinite_loss}
Let $L$ be as in Lemma \ref{continuous_version}, $\bar G_0$ be as in Assumption \ref{A2}, 
$\lambda_t$ be as in (\ref{lambda}), and let $x \in \mathbb{R}$. As $|h| \to -\infty$ we have that
\begin{align}\label{ruin}
\sup_{j \in \mathbb{R}^{t-1}}L(x,j,h) \to -\infty
\end{align}
almost surely. (In the case $t=1$ we mean that $L$ does not depend on $j$ and there is no supremum.)
\end{lemma}
\begin{proof}
Without loss of generality we assume $t=2$, the case of $t>2$ being only notationally more cumbersome. 
By Assumption \ref{A2}, the definition of $L$ and using (\ref{lambda}), we have for all $x,j,h \in \mathbb{R}$ that
\begin{align*}
L(x,j,h) &= E[V(x,j,h) | \mathcal{H}]
\\ 
&\leq E[\bar G_0(x + \kappa_{2}(j,h))| \mathcal{H}]
\\
&
\leq 
E[\bar G_0(x + \lambda_{2}(h))| \mathcal{H}]
\end{align*}
almost surely. Apply Fatou's reverse lemma to the inequalities above. 
Considering (\ref{lambda}) and the absence of dependence on the variable $j$, continuity of $\bar G_0$ shows (\ref{ruin}).
\end{proof}

\begin{lemma}\label{lemma::monotonicity_}
The inequality
\begin{align}\label{pók}
\begin{split}
&L(x_1, j, h) \leq L(x_2, j, h),
\end{split}
\end{align}
holds for all $x_1,x_2, h \in  \mathbb{R}$ with $x_1 < x_2$, for all $j \in \mathbb{R}^{t-1}$, and for almost every $\omega \in \Omega$.
\end{lemma}

\begin{proof} Without loss of generality we can assume $t = 2$. let $\Omega'$ be a $\mathbb{P}$-full measure set such that \eqref{pók} holds for all $x_{1},x_{2},j,h\in\mathbb{Q}$
on $\Omega'$. Let $\Omega''$ be the full measure set on which $L$ is continuous. Then on $\Omega'\cap\Omega''$ \eqref{pók} holds
for \emph{all} $x_{1},x_{2},j,h\in\mathbb{R}$, by continuity.	

% Let us denote with $\bar \Omega$ the full measure set where the assertion of Lemma \ref{continuous_version} holds. Furthermore, let us introduce the set
%$$I = \{ (x_1,x_2,j,h)  \ : \ x_1, x_2,h \in \mathbb{R}, x_1<x_2, j \in \mathbb{R}  \}$$ and for any $z = (x_1, x_2, j, h) \in I$, let us denote with $A_{z}$ the subset of $\bar \Omega$ for which the inequality
%\begin{align*}
%\begin{split}
%L(x_1, j, h) \leq L(x_2, j, h),
%\end{split}
%\end{align*}
%holds. Also, employ the notations $R = \cup_{z \in I} (\bar \Omega - A_z)$, and $Q = \cup_{q \in \mathbb{Q}^{2+t} \cup I} (\bar\Omega - A_q)$. Let $\omega_0 \in R$. Then, by definition of $R$, there exists $z^{(0)} = (x_1^{(0)}, x_2^{(0)}, j^{(0)}, h^{(0)}) \in I$ so that $\omega_0 \in \bar \Omega - A_{z^{0}}$. Continuity of $L(\omega_0, \cdot, \cdot, \cdot)$, implies that there exist  $\varepsilon > 0$, and a sequence $q(k) = (x_1^{(0)}(k), x_2^{(0)}(k), j^{(0)}(k), h^{(0)}(k)) , \ k\in \mathbb{N}$ in $\mathbb{Q}^{2 + t} \cap I$, so that $q(k) \to z^{(0)}$, and we also have
%\begin{align}
%\begin{split}
%L(\omega_0, x_1^{(0)}(k), &j^{(0)}(k), h^{(0)}(k)) -\varepsilon = L(\omega_0, x_1^{(0)}, j^{(0)}, h^{(0)}) -\varepsilon
%\\
%&> 
%L(\omega_0, x_2^{(0)}, j^{(0)}, h^{(0)}) = L(\omega_0, x_2^{(0)}(k), j^{(0)}(k), h^{(0)}(k)).
%\end{split}
%\end{align}
%This means that there exists an index $k_0 \in \mathbb{N}$ so that $\omega_0 \in \bar \Omega - A_{q(k_0)}$, meaning that $\omega_0 \in Q$. Since $\omega_0$ was arbitrary, we have proved that $R \subset Q$. Note, that as $L(x,j,h)$ is a version of $E[V(x,j,h)|\mathcal{H}]$ for any $x,j,h$, and $V$, in its first variable, inherits the non-decreasing property from $G$, we have that for all $q \in \mathbb{Q}^{2+t}\cup I$ the event $A_q$ is a full measure set. This, in return, implies -- referencing the completeness of the sigma field $\mathcal{F}$ -- that $R \subset \Omega$ is, in fact, an event, and it is one with zero measure.
\end{proof}

\begin{lemma}\label{lemma::kettes}
Fix $l \in \mathbb{Z}$ and $m \in \mathbb{N}$. There exits an $\mathcal{H}$-measurable $K = K(l,m)$ such that 
\begin{equation}\label{majus6}
L(x,j, h) \leq L(x, j, h \indicator{|h|\leq K}),
\end{equation}
for all $x \in [l,l+1]$, $j \in [-m,m]^{t-1} \subset \mathbb{R}^{t-1}$, $h \in \mathbb{R}$, and for almost every $\omega \in \Omega$.
\end{lemma}

\begin{proof}
Without loss of generality assume $t=2$, and let $\bar \Omega$ denote the full measure set 
where the conclusion of Lemma \ref{continuous_version} and Lemma \ref{lemma::monotonicity_} hold. 
For every $\omega \in \bar \Omega$ choose a (measurable) $m^{+}(h) = m^{+}(\omega, l,h)$ such that  
\begin{align*}
L(l+1, j, h)(\omega) \leq 
L(l+1, m^{+}(h), h)(\omega)
\end{align*}
holds for all $j \in [-m,m]$ and $h \in \mathbb{R}$, this is possible by continuity of $L$. Likewise, for every $\omega \in\bar \Omega$ choose $m^{-} = m^{-}(\omega, l)$ such that  
\begin{align*}
L(l, m^{-}, 0)(\omega) \leq 
L(l, j, 0)(\omega)
\end{align*}
holds for all $j\in[-m,m]$.
Fix $l \in \mathbb{Z}$ and $m \in \mathbb{N}$. Using Lemma \ref{lemma::infinite_loss} for all $\omega \in \bar\Omega$ there exists $K(\omega) = K(\omega, l, m)$ so that for all $h \in \mathbb{R}$ it holds that
\begin{align}\label{M}
\begin{split}
|h| > K(\omega) \implies L(l+1, m^{+}(h), h)(\omega) \leq L(l, m^{-}, 0)(\omega).
\end{split}
\end{align}
Now note, that $\omega \to K(\omega)$ can be chosen in a way that it is $\mathcal{H}$-measurable as a random variable. On the event $\{|h| > K\} \cap \bar \Omega$, using Lemma \ref{lemma::monotonicity_}, and the statement in (\ref{M}), we have 
\begin{align}\label{boundidstratineq}
\begin{split}
&L(x, j, h) \leq L(l+1, j, h) \leq 
L(l+1, m^{+}(h), h),
\\
&\leq L(l, m^{-}, 0) \leq  L(l, j, 0) \leq 
L(x, j, 0),
\end{split}
\end{align}
for every $x \in [l,l+1]$, $j \in [-m,m]$, $h\in \mathbb{R}$, and for all $\omega \in \bar \Omega$: completing the argument.
\end{proof}


\begin{lemma}\label{lemma::pointwise_essup}
There exists an $\mathcal{H}\otimes\mathcal{B}(\mathbb{R})$-measurable function 
$$G :\Omega\times\mathbb{R}\times \mathbb{R}^{t-1} \mapsto \mathbb{R},$$ such that $G$ is 
continuous almost surely, $G$ is non-decreasing in its first variable for almost all 
$\omega\in\Omega$, furthermore, for all $x$, and for all $j \in \mathbb{R}^{t-1}$ we have 
\begin{align}\label{majus5}
\begin{split}
G(x, j) =  \mathrm{ess.}\sup_{H \in \Xi} L(x, j, H),
\end{split}
\end{align}
almost surely.
\end{lemma}

\begin{proof} We follow arguments of Lemma 3.17 in \cite{laurence}.
Let $x,j \in \mathbb{R}$, and without loss of generality assume $t = 2$ and set $l\in\mathbb{Z}$ 
and $m \in \mathbb{N}$ so that $x \in [l,l+1]$ and $j \in [-m,m]$ holds. We will 
work out the statement in consideration elementwise on $\Omega$, and to this end -- out of the usual -- we 
do not omit to display dependence on $\omega\in\Omega$ throughout the proof, until further notice. 
Denote with $\bar \Omega$ the full measure set for which the conclusions of Lemma \ref{continuous_version}, 
\ref{lemma::monotonicity_}, \ref{lemma::kettes} hold, and let $\omega \in \bar \Omega$. Let us define  
\begin{align*}
 B(\omega, x, j) = \sup_{h \in \mathbb{Q}, |h| \leq K(\omega)}L(\omega, x,j,h) = \sup_{h \in \mathbb{R}, |h| \leq K(\omega)}L(\omega, x,j,h),
\end{align*}
where $K(\omega) = K(\omega, l,m) = K(l,m)$ is as in Lemma \ref{lemma::kettes}. 
This is measurable, being the supremum of countably many measurable functions. Let us fix furthermore a sequence 
$(x_n,j_n)_{n \in \mathbb{N}} \subset ([l,l+1]\times[-m,m])\cap \mathbb{R}^2$ for which $(x_n,j_n) \to (x,j)$. 
Observe, that by definition of $B$, for every $k \in \mathbb{N}$ there exists $h_k(\omega, x,j)$, with 
$h_k(\omega, x,j) \leq K(\omega)$, so that $B(\omega, x,j) -1/k \leq L(\omega, x,j,h_k(\omega, x,j))$. 
The fact that for all $n\in\mathbb{N}$ we have $B(\omega,x_n,j_n) \geq L(\omega,x_n,j_n,h_k(\omega, x,j))$, along with continuity of $L$ yields
\begin{align*}
\liminf_{n}B(\omega,x_n,j_n) \geq & L(\omega,x,j,h_k(\omega, x,j)) 
\\
&
\geq
B(\omega, x,j) -1/k,
\end{align*}
which in the limiting case of $k \to \infty$ means $\liminf_{n}B(\omega,x_n,j_n) \geq B(\omega, x,j)$. 

Take a sequence $\{n_k, \ k\in\mathbb{N}\} \subset \mathbb{N}$ so that
\begin{align}\label{moha}
\limsup_{n}B(\omega,x_n,j_n) = \lim_{k}B(\omega,x_{n_k},j_{n_k}).
\end{align}
Since $\{h : h\in\mathbb{Q}, \ |h|\leq K(\omega) \}$ is a precompact set in $\mathbb{R}$, for every 
$k\in\mathbb{N}$ there exists $\mathbb{R} \ni h_{n_k}^{*}(\omega) \leq K(\omega)$ so that 
$B(\omega,x_{n_k},j_{n_k}) = L(\omega,x_{n_k},j_{n_k},h_{n_k}^{*}(\omega))$. Using  the compactness of the 
closure there exists $\mathbb{R}\ni h^*(\omega) \leq K(\omega)$ and a subsequence 
$(a_k)_{k\in\mathbb{N}}$ of $\{n_k : k\in\mathbb{N}\}$ so 
that $h^{*}_{a_k}(\omega) \to h^*(\omega)$, $k\to\infty$. These, and (\ref{moha}) imply
\begin{align*}
\limsup_{n} & B(\omega,x_n,j_n) = \lim_{k}B(\omega,x_{n_k},j_{n_k}) = \lim_{k}B(\omega,x_{a_k},j_{a_k})
\\
&
= \lim_{k}L(\omega,x_{a_k},j_{a_k},h_{a_k}^{*}(\omega)) = L(\omega,x,j,h^{*}(\omega)) \leq B(\omega,x,j),
\end{align*}
establishing the continuity of $B$. 

As far as monotonicity is concerned, the mapping $x \to B(\omega,x,j)$ inherits the 
non-decreasing property from $L$ (stated in Lemma \ref{lemma::monotonicity_}) naturally.

From Lemma \ref{gast} it follows that
\begin{align*}
B(\omega,x,j) = \mbox{ess.sup}_{H(\omega) \leq K(\omega)} L(\omega, x,j, H(\omega)).
\end{align*}
In the discussion above $\omega \in \bar \Omega$ was arbitrary, and returning to the usual practice of not displaying the dependence on it, Lemma \ref{lemma::kettes} shows that 
\begin{align*}
\mbox{ess.sup}_{H \in \Xi} L( x,j, H)  \leq \mbox{ess.sup}_{\Xi \ni H \leq K}& L(x,j, H) = B(x,j) 
 \\
 &\leq
 \mbox{ess.sup}_{H \in \Xi} L( x,j, H)
\end{align*}
holds almost surely, finishing the argument.

\end{proof}

%L(\omega,x,j,h_n^{*}(\omega))
%\begin{lemma}\label{lemma::pointwise_essup}
%There exists an $\mathcal{H}\otimes\mathcal{B}(\mathbb{R})$-measurable function $$G :\Omega\times\mathbb{R}\times \mathbb{R}^{t-1} \mapsto \mathbb{R},$$ such that for
%almost all $\omega\in\Omega$, $G(\cdot)$ is right-continuous and non-decreasing in its first variable, left continuous and non-increasing in the rest \textcolor{red}{different cases for the parameters} of its real variables, and for all $x$, and for all $j \in \mathbb{R}^{t-1}$ we have 
%\begin{align}\label{majus5}
%\begin{split}
%G(x, j) =  \mathrm{ess.}\sup_{H \in \Xi} L(x, j, H),
%\end{split}
%\end{align}
%almost surely.
%\end{lemma}
%
%\begin{proof}
%We will assume, without loss of generality, that $t=2$. Let $x, j \in \mathbb{R}$, and for each of these choose a jointly measurable representative of the essential supremum in (\ref{majus5}), and denote the version with $F(x,j)$. We have, by Assumption \ref{A1}, that there exists $A^{-}, \ A^{+}  \in \mathcal{G}$ with $A^{-}\cup A^{+}$ being an event of probability $1$, and a $\mathcal{G}$-measurable $\varepsilon$ so that on the random rational region $\mathbb{Q}\times \left( [j-\varepsilon(\omega),j+\varepsilon(\omega)]\cap\mathbb{Q} \right)$ the mapping $j \to F(x,j)$ is non-decreasing when $\omega \in A^{-}$ and non-increasing when $\omega \in A^{+}$. Without loss of generality we assume that $\omega \in A^{+}$. Now, let us define
%\begin{equation}
%v(x,j)=
%\inf_{j - \varepsilon < b < j, b \in \mathbb{Q}} \inf_{a > x, a \in \mathbb{Q}}  F(a,b).
%\end{equation}
%The function $(x,j) \to v(x,j)$ is $\mathcal{H} \otimes \mathcal{B}(\mathbb{R}^2)$-measurable, since the infima is taken over countable sets (and one can express $v(x,j)$ with the help of the bookkeeping functions $f_{a,b}(x,j,\omega) = \indicator{x<a}\indicator{j<b}F(a,b,\omega)$). 
%
% The function $v$ is right semi-continuous in its first variable and left semi-continuous in its second variable by construction, and it is inherits monotonicity properties of $F$ by definition.
% 
%Obviously, for any $x,j$ we have $F(x,j) \leq v(x,j)$ almost surely. To see the reverse inequality we give the following argument.
%
%Fix $x,j \in \mathbb{R}$, and choose $l \in \mathbb{Z}$ and $m \in \mathbb{N}$ so that $x \in [l, l+1]$ and $j \in [-m,m]$. Choose furthermore a sequence $\mathbb{Q}^2\supset \{(q_n,p_n), \ n\in \mathbb{N}\}$, so that for all $n \in \mathbb{N}$ we have  $x < q_n, j>p_n$, $(q_n,p_n) \to (x,j)$. Using Lemma \ref{lemma::kettes}, and the definition of $F$, we have that there exists $\mathcal{H}$-measurable $K = K(l,m)$, and for all $n \in \mathbb{N}$ there exists $\mathcal{H}$-measurable $H_n$, so that 
%\begin{align*}
%\begin{split}
%F(q_n,p_n)& - 1/n \leq L(x,j, H_n) \leq L(x, j, H_n \indicator{|H_n| \leq K}).
%\end{split}
%\end{align*}
%Using Lemma 2 of \cite{KabanovStricker2001}, there exists $\mathcal{H}$-measurable $\bar H$ and a sequence of $\mathcal{H}$-measurable indices $\{ k_n, \ n \in \mathbb{N} \} \subset \Xi(\mathcal{H})$, so that the subsequence $H_{k_n} \indicator{|H_{k_n}| \leq K}, \ n \in \mathbb{N}$ converges almost surely to $\bar H$, as $n \to \infty$. Using this and continuity of $L$, we acquire
%\begin{align*}
%\begin{split}
%v(x,j) &= \lim_{n \to \infty} \Big( F(q_{k_n}, p_{k_n}) - 1/k_n \Big) \leq \lim_{n \to \infty} L(x,j, H_{k_n} \indicator{|H_{k_n}| \leq K})
%\\
%=&
%L(x,j, \bar H \indicator{|\bar H| \leq K}) \leq F(x,j).
%\end{split}
%\end{align*}  
%This claim shows that $v$, for every $x,j \in \mathbb{R}$, is indeed a version of the essential supremum in consideration, and the choice $G = v$ finishes the argument.
%
%\end{proof}

\begin{lemma}\label{lemma::harmas}
Let $X,H_1,\hdots, H_{t-1} $ be $\mathcal{H}$-measurable random variables, and let $G$ be as in Lemma \ref{lemma::kettes}. Then, the quantity $G(X,(H_1,\hdots, H_{t-1}))$ is a version of the essential supremum
\begin{align*}
\begin{split}
\mathrm{ess.}\sup_{H\in \Xi_{t-1}}
L(X,(H_1,\hdots,H_{t-1}), H).
\end{split}
\end{align*}
\end{lemma}

\begin{proof}
Without loss of generality we assume that $t = 2$, and that $X$ and $H_1$ take values in $[l,l+1]$ and $[-m,m]$ 
respectively. Let $X_n$ and $H_1^{(n)}$ be $\mathcal{H}$-measurable random variables, taking values in $[l,l+1]\cap\mathbb{Q}$ and 
$[-m,m]\cap\mathbb{Q}$, respectively, for all $n \in \mathbb{N}$, and possessing also 
the limiting properties $X_n \to X$ and $H_1^{(n)} \to H_1$.

Observe that for all $n\in\mathbb{N}$, on a full measure set, we have that $$\mathrm{ess.sup}_{H\in \Xi_{t-1}} L(X_n,H_1^{(n)},H) = G(X_n,H_1^{(n)}),$$ and furthermore, as a consequence, for all $n\in\mathbb{N}$ there exists $H_n$ so that 
\begin{align}\label{fin1}
 L(X_n,H_1^{(n)},H_n) \geq  G(X_n,H_1^{(n)}) - 1/n
\end{align}
almost surely. According to Lemma \ref{lemma::kettes} there exists an $\mathcal{H}$-measurable $K = K(l,m) = K(\omega,l,m)$ so that
\begin{align}\label{fin2}
 L(X_n,H_1^{(n)},H_n \indicator{|H_n|\leq K}) \geq  L(X_n,H_1^{(n)},H_n).
\end{align}
Note that the $K$ does not depend on the integer $n$ in any way. Putting together (\ref{fin1}) and (\ref{fin2}) gives, for all $n \in \mathbb{N}$ the almost sure inequality
\begin{align}\label{fin3}
L(X_n,H_1^{(n)},H_n \indicator{|H_n|\leq K}) \geq  G(X_n,H_1^{(n)}) - 1/n.
\end{align}
Now using Lemma \ref{randomsub}, there exists an $\mathcal{H}$-measurable subsequence $k_n = k_n(\omega), \ k\in\mathbb{N}$ and an $\mathcal{H}$-measurable $\bar H$ so that $H_{k_n} \indicator{|H_{k_n}|\leq K} \to \bar H$ holds almost surely. In (\ref{fin3}) taking the limit as $n \to \infty$ along the sequence $k_n, \ n\in\mathbb{N}$, and utilizing continuity of $L$ and $G$ yields
\begin{align*}
L(X,H_1,\bar H) \geq G(X,H_1),
\end{align*}
which in return implies $\mathrm{ess.sup}_{H\in \Xi_{t-1}}L(X,H_1,H) \geq G(X,H_1)$.

On the other hand, by definition of $G$, for every $\mathcal{H}$-measurable $H$ we have 
\begin{align*}
L(X,H_1,H) =& \lim_{n}L(X_n,H_1^{(n)},H) \leq \lim_{n} \mbox{ess.sup}_{H}L(X_n,H_1^{(n)},H)
\\
&=
 \lim_{n} G(X_n, H_1^{(n)}) = G(X,H_1).
\end{align*}
Taking the essential supremum of both sides yields the inequality
\begin{align*}
\mathrm{ess.sup}_{H\in \Xi_{t-1}}L(X,H_1,H) \leq G(X, H_1)
\end{align*}
on a full measure set: finishing the proof.
\end{proof}

%\begin{proof}
%Without the loss of generality assume that $t = 2$, $X$ takes values in $[l,l+1]$ for some $l\in \mathbb{Z}$, and $H_1$ takes values in $[-m,m]$ for some $m\in \mathbb{Z}$. Let $X_n, n \in \mathbb{N}$ and  $H_1^{(n)}, n \in \mathbb{N}$ be a sequence of countable valued random variables, taking values in $[l,l+1]$ and $[-m,m]$ respectively, so that $X_n \to X, \ H_1^{(n)} \to H_1$, in the almost sure sense. Using Lemma \ref{lemma::pointwise_essup} let $G$ be as in (\ref{majus5}), and observe that its properties imply
%\begin{equation}\label{fotel}
%G(X_n, H_1^{(n)}) \to G(X, H_1)
%\end{equation}
%as $n \to \infty$, almost surely. Furthermore, using Lemma \ref{lemma::monotonicity_}, for all $\mathcal{H}$-measurable $H$ and for all $n \in \mathbb{N}$ we have
%\begin{equation}\label{asztal}
%L(X,H_1, H) \leq L(X_n,H_1 ,H)
%\end{equation}
%almost surely. Note that the statement that this proof aims to establish (that of Lemma \ref{lemma::harmas}) trivially holds for countable valued random variables. Using the latter fact, in (\ref{asztal}), take the essential supremum of both sides, and then, considering (\ref{fotel}), pass the expressions to the limit to see that
%\begin{align*}
%\begin{split}
%\mathrm{ess.}\sup_{H\in \Xi_{t-1}}
%L(X, H_1,H) \leq G(X,H_1)
%\end{split}
%\end{align*}
%almost surely. 
%
%Now we prove the reverse inequality. Using the construction of $G$, Lemma \ref{lemma::kettes}, and the fact that the $X_n, \ H_n$ are countable valued, there exists $K = K(l,m)$, and for every $n \in \mathbb{N}$ there exists an $\mathcal{H}$-measurable $H_n$ so that
%\begin{equation}
%G(X_n,H_1^{(n)}) - \frac{1}{n} \leq L(X_n,H_1^{(n)}, H_n) \leq L(X_n, H_1^{(n)}, H_n \indicator{|H_n| \leq K}).
%\end{equation}
%Utilizing Lemma X of REF, there exists an $\mathcal{H}$-measurable variable $\bar H$, and an $\mathcal{H}$-measurable index set $n_k, \ k \in \mathbb{N}$ so that -- using the notation $\bar H_n = H_n \indicator{|H_n| \leq K}$ -- we have $\bar H_n \to \bar H$. Utilizing continuity of $L$, and passing to limits yield
%\begin{equation}
%G(X,H_1) \leq L(X, H_1, \bar H \indicator{|\bar H| \leq K})
%\end{equation}
%almost surely: finishing the proof.
%\end{proof}

%\begin{lemma}
%Let $G$ be as in Lemma \ref{lemma::pointwise_essup}. Outside a negligible set the function $(x,j) \to G(x,j)$ is continuous.
%\end{lemma}
%
%\begin{proof}
%Let us assume without loss of generality that $t =2$ and $h<0$ \textcolor{red}{wlog jo lesz?}. Let $\varepsilon>0$ and consider the set 
%\begin{align*}
%\mathcal{E} = \{ (\omega, x,j) \ : \ G(\omega, x,j) - \varepsilon \geq \sup_{j < b, b \in \mathbb{Q}} \sup_{a < x, a \in \mathbb{Q}} G(\omega, a,b)  \},
%\end{align*}
%and assume that its projection on $\Omega$, which we denote by $A \in \mathcal{H}$, is of positive probability. Denote with $(X, J)$ an $\mathcal{H}$-measurable selector of $\mathcal{E}$ on $A$. Let furthermore $X_k, \ J_k$ be $\mathcal{H}$-measurable rational valued functions so that $(X_k,J_k) \to (X,J)$ and $X_k < X$ and $J_k > J$.  Choose for each $l \in \mathbb{N}$ an $\mathcal{H}$-measurable $H_l$ so that $G(X,J) - 1/l \leq L(X,J,H_l)$. On the event $A$ for every $l \in \mathbb{N}$ we have
%\begin{align*}
%L(X_k,J_k, H_l) \leq G(X_k,J_k) \leq G(X,J) - \varepsilon.
%\end{align*}
%Since $L$ is continuous outside a negligible set, taking the limit  $k \to \infty$ yields 
%$$G(X,J) - 1/l \leq G(X,J) - \varepsilon,$$
%a contradiction in the limiting case $l \to \infty$.
%\end{proof}

\begin{lemma}\label{optimal_H}
Let $G$ be as in Lemma \ref{lemma::pointwise_essup}. Let $X, H_1,\hdots, H_{t-1}$ be $\mathcal{H}$-measurable random variables. Then, there exists an $\mathcal{H}$-measurable $H^*$ so that
\begin{align*}
G(X,(H_1,\hdots,H_{t-1})) = L(X,(H_1,\hdots, H_{t-1}), H^*)
\end{align*}
holds almost surely.
\end{lemma}

\begin{proof}
Without loss of generality assume that $t=2$, $X$ almost surely takes values in the closed interval $[l,l+1]$, $H_1$ takes values in the closed interval $[-m,m]$ for some $l \in \mathbb{Z}$ and for some $m \in \mathbb{N}$. Let us define
\begin{equation*}
\mathcal{O} = \Big\{  L(X,H_1, H), \ \ H \in \Xi(\mathcal{H})   \Big\}.
\end{equation*}
This set is upward directed in terms of the almost sure sense of the usual "less than or equal" relation. Thus, using Proposition VI.1.1. of \cite{neveu} we have that there exists a sequence $\{ H_n, \ n\in\mathbb{N} \} \subset \Xi(\mathcal{H})$ for which the limiting property
\begin{align}\label{krokodil0}
\begin{split}
 L(X, H_1, H_n ) \to G(X,H_1) 
\end{split}
\end{align}
holds almost surely as $n \to \infty$. Utilizing Lemma \ref{lemma::kettes}, and the $K = K(l,m)$ within, we have for each $n \in\mathbb{N}$ that
\begin{align}\label{beka0}
\begin{split}
L&(  X, H_1, H_n) \leq L(X, H_1, H_n \indicator{|H_n| \leq K}),
\end{split}
\end{align}
almost surely. Then, again using Lemma \ref{randomsub}, there exists an $\mathcal{H}$-measurable subsequence, 
say $k_n, \ n \in \mathbb{N}$, and there exists an $\mathcal{H}$-measurable $H^*$ so that 
$$
H_{k_n} \indicator{|H_{k_n}| \leq K} \to H^*
$$
in the almost sure sense. Continuity of $L$, (\ref{krokodil0}), and (\ref{beka0}) together guarantees
\begin{align*}
\begin{split}
G(X,H_1) &= \lim_{n \to \infty} L(X,H_1, H_{k_n} ) 
\\
\leq&
\lim_{n \to \infty} L(X, H_{k_n} \indicator{|H_{k_n}| \leq K}) = L(X,H_1, H^*).
\end{split}
\end{align*}
The proof is complete.
\end{proof}


\begin{lemma}\label{x_bound}
Let $G$ be as in Lemma \ref{lemma::pointwise_essup}. There exists $\bar G : \mathbb{R} \to \mathbb{R}$, and $C>0$ such that the 
following requirements are met: as $x \to -\infty$ we have
$$\bar G(x) \to - \infty,$$
almost surely and, for all $x \in \mathbb{R}$, $j \in \mathbb{R}^{t-1}$ we have 
\begin{equation*}
G(x,j) \leq \bar G(x) \leq C
\end{equation*}
in the almost sure sense. 
%Furthermore, $G(x,j) \leq C$ for all $(x,j) \in \mathbb{R}\times\mathbb{R}^{t-1}\times\mathbb{R}$ and for almost every $\omega\in\Omega$.
\end{lemma}
\begin{proof}
Let $G$ and $H^*$ be as in Lemma \ref{lemma::pointwise_essup} and Lemma \ref{optimal_H} respectively. Without loss of generality we 
can assume $t=2$. Let $H^*$ be as in Lemma \ref{optimal_H}, and let us choose a sequence of rational-valued random variables $H_n^*$ increasing to it. Using the market bound for the innovations $\kappa_{\cdot}$ in (\ref{lambda}), Assumption \ref{A3}, and Fatous'  reverse Lemma, we have for $x,j \in \mathbb{R}$  that
\begin{align*}
G(x,j) =& L(x,j,H^*) =  \lim_{n} L(x,j,H^*_n) \leq \limsup_{n} E[V(x,j,H_n^{*})|\mathcal{H}] 
\\
&\leq \limsup_{n} E[\bar G_0(x +\kappa_2(j, H^*_n))|\mathcal{H}] \leq E[\bar G_0(x +\kappa_2(j, H^*))|\mathcal{H}]
\\
&
\leq
E[\bar G_0(x + P_t^2 \delta_t / 4)| \mathcal{H}] =:G^{'}(x)
\end{align*}
almost surely. 

Fix some $v_0 \in \mathbb{R}$ and note that we have $G_0(x,v_0) \leq \bar{G}_0(x)$. With Assumption \ref{A3} in mind, 
take a continuous version of $G^{'}$ using Lemma \ref{szeder}, and denote it by $G^{''}$. By construction for a.e. $\omega\in \Omega$ we have $G(\cdot,\cdot) \leq C$ and $G^{''}(\cdot) < C$, and observe that by the reverse Fatou lemma we have $G^{''}(x) \to -\infty$ and $x \to -\infty$. The function $G^{''}(\cdot)$ inherits monotonicity from $\bar G_0$. The former fact can be seen with similar reasoning given in the proof of Lemma \ref{lemma::monotonicity_}. The choice $\bar G = G^{''}$ completes the proof.
\end{proof}

%\begin{lemma}
%Let $G$ and  $\bar G$ be as in Lemma \ref{lemma::pointwise_essup} and Lemma \ref{x_bound} respectively, and consider the function
%\begin{align*}
%V_1:\Omega\times & \mathbb{R} \times \mathbb{R}^{t-2} \times \mathbb{R} \to \mathbb{R}, 
%\\
%&
%\ (x,j,h) \mapsto G(x + \kappa_{t-1}(j_1, \hdots, j_{t-2},h),(j_1, \hdots, j_{t-2},h)).
%\end{align*}
%For any $m\in \mathbb{N}$, setting $M = [-m,m]^{t}$,  there exists an integrable random variable $\hat M = \hat M(m)$ so that $\hat M \leq V_1(x,j,h)$ and $\hat M \leq \bar G(x)$ holds for every $(x,j,h) \in M$ and for almost every $\omega \in \Omega$.
%\end{lemma}
%
%\begin{proof}
%Let $m \in \mathbb{N}$. Using Assumption \ref{A3} there exists and integrable, $\mathcal{G}$-measurable $M_1 = M_1(m,t)$ so that 
%\begin{align*}
%G_0(x + \kappa_{t-1}(j), (j, 0)) \geq M_1
%\end{align*}
%for all $x \in [-m,m]$, for all $j \in [-m,m]^{t-1}$, and for almost every $\omega\in\Omega$. For fixed $j = (j_1, \hdots, j_{t-1}) \in \mathbb{R}^{t-1} \cap [-m,m]^{t-1}$ and $x \in [-m,m]$, with the notation $(j,0) = (j_1,\hdots, j_{t-1}, 0)$, we have  
%\begin{align}\label{hamv}
%\begin{split}
%G(x + &\kappa_{t-1}(j),j) \geq L(x + \kappa_{t-1}(j),j,0)
%\\
%&
%= E[V(x + \kappa_{t-1}(j), j, 0)| \mathcal{H}] 
%\\
%&
%= E[G_0(x + \kappa_{t-1}(j) + \kappa_{t}((j, 0)), (j, 0))| \mathcal{H}]
%\\
%&
%= E[G_0(x + \kappa_{t-1}(j), (j, 0))| \mathcal{H}] = E[M_1| \mathcal{H}]
%\end{split}
%\end{align}
%almost surely. Since $G$ is continuous the relation established in (\ref{hamv}) above holds for all $x \in [-m,m]$, for all $j \in [-m,m]^{t-1}$, and for almost every $\omega\in\Omega$.
%
%Now we turn our attention to consturcting the lower bound for $\bar G_0$. Observe that for all $x \in \mathbb{R}$ we have by construction of $\bar G$ that
%\begin{align*}
%\bar G(x) &\geq G(x,0) \geq L(x,0,0) = E[V(x,0, 0) |\mathcal{H}] = E[G_0(x, 0, 0) |\mathcal{H}]
%\end{align*}
%almost surely. Thus, using (\ref{hamv}) with $j = 0$ yields $\bar G(x) \geq E[M_1| \mathcal{H}]$ almost surely, from which by a standard argument, again utilizing the continuity of $\bar G$ gives $\bar G (x) \geq E[M_1| \mathcal{H}]$ for all $x \in \mathbb{R}$, and almost every $\omega\in\Omega$. The choice $\hat M = E[M_1| \mathcal{H}]$ gives the desired lower bound and finishes the proof.
%\end{proof}

%\begin{lemma}
%Let $G$ and $\bar G$ be as in Lemma \ref{lemma::pointwise_essup} and Lemma \ref{x_bound} respectively, and let $t^{'} \in \{1,\hdots,t-1\}$. For any $m\in \mathbb{N}$ there exists an integrable random variable $ M_1 = M_1(m,t^{'})$ so that $$M_1 \leq G(x + \kappa_{t^{'}}((j_1,\hdots, j_{t^{'}})),j)$$ for every $x \in \mathbb{R}$, $j \in [-m,m]^{t-1}$, and  for almost every $\omega \in \Omega$. Furthermore, there exists an integrable random variable $\bar M_1 = \bar M_1(m, t^{'})$ so that
%$$\bar M_1 \leq \bar G(x + P_{t^{'}}^2 \delta_{t^{'}} / 4)$$
%holds for every $x  \in [-m,m]$ and for almost every $\omega \in \Omega$.
%\end{lemma}
%
%\begin{proof}
%Let $m \in \mathbb{N}$ and $t^{'} \in \{1,\hdots,t-1\} $. Using Assumption \ref{A3} there exists and integrable, $\mathcal{G}$-measurable $M_0 = M_0(m,t^{'})$ so that with with the notation $(j,0) = (j_1,\hdots, j_{t-1}, 0)$ we have 
%\begin{align*}
%G_0(x + \kappa_{t^{'}}((j_1,\hdots,j_{t^{'}})), (j, 0)) \geq M_0
%\end{align*}
%for all $x \in [-m,m]$, for all $j \in [-m,m]^{t-1}$, and for almost every $\omega\in\Omega$. For fixed $j = (j_1, \hdots, j_{t-1}) \in [-m,m]^{t-1}$ and $x \in [-m,m]$, we have  
%\begin{align}\label{hamv}
%\begin{split}
%G(x + &\kappa_{t^{'}}((j_1,\hdots,j_{t^{'}})),j) \geq L(x + \kappa_{t^{'}}((j_1,\hdots,j_{t^{'}})),j,0)
%\\
%&
%= E[V(x + \kappa_{t^{'}}((j_1,\hdots,j_{t^{'}})), j, 0)| \mathcal{H}] 
%\\
%&
%= E[G_0(x + \kappa_{t^{'}}((j_1,\hdots,j_{t^{'}})) + \kappa_{t}((j, 0)), (j, 0))| \mathcal{H}]
%\\
%&
%= E[G_0(x + \kappa_{t^{'}}((j_1,\hdots,j_{t^{'}})), (j, 0))| \mathcal{H}] = E[M_0| \mathcal{H}]
%\end{split}
%\end{align}
%almost surely. Since $G$ is continuous the relation established in (\ref{hamv}) above holds for all $x \in [-m,m]$, for all $j \in [-m,m]^{t-1}$, and for almost every $\omega\in\Omega$. The choice $M_1 = E[M_0| \mathcal{H}]$ gives a desired lower bound. We now turn our attention to constructing the bound for $\bar G$. To this end using Assumption \ref{A3} there exists $\bar M = \bar M(m, t^{'})$ so that for $x \in [-m,m]$ we have
%\begin{equation}\label{julikutya}
%\bar G(x) = E[\bar G_0(x + P_{t^{'}}^2 \delta_{t^{'}} / 4)| \mathcal{H}] \geq E[\bar M| \mathcal{H}]
%\end{equation}
%almost surely. Using continuity of $\bar G$ we have that (\ref{julikutya}) holds for all $x \in [-m,m]$ and for almost every $\omega\in\Omega$, and setting $\bar M_1 = E[\bar M| \mathcal{H}]$ finishes the proof.
%\end{proof}


\begin{lemma}\label{lemma::lower_bd}
Let $G$ be as in Lemma \ref{lemma::pointwise_essup}, and let $t^{'} \in \{1,\hdots,t-1\}$. For any $m\in \mathbb{N}$ there exists an integrable random variable $ M_1 = M_1(m,t^{'})$ so that $$M_1 \leq G(x + \kappa_{t^{'}}((j_1,\hdots, j_{t^{'}})),j)$$ for every $x \in \mathbb{R}$, $j \in [-m,m]^{t-1}$, and  for almost every $\omega \in \Omega$.
\end{lemma}

\begin{proof}
Let $m \in \mathbb{N}$ and $t^{'} \in \{1,\hdots,t-1\} $. Using Assumption \ref{A3} there exists and integrable, $\mathcal{G}$-measurable $M_0 = M_0(m,t^{'})$ so that with with the notation $(j,0) = (j_1,\hdots, j_{t-1}, 0)$ we have 
\begin{align*}
G_0(x + \kappa_{t^{'}}((j_1,\hdots,j_{t^{'}})), (j, 0)) \geq M_0
\end{align*}
for all $x \in [-m,m]$, for all $j \in [-m,m]^{t-1}$, and for almost every $\omega\in\Omega$. For fixed $j = (j_1, \hdots, j_{t-1}) \in [-m,m]^{t-1}$ and $x \in [-m,m]$, we have  
\begin{align}\label{hamv}
\begin{split}
G(x + &\kappa_{t^{'}}((j_1,\hdots,j_{t^{'}})),j) \geq L(x + \kappa_{t^{'}}((j_1,\hdots,j_{t^{'}})),j,0)
\\
&
= E[V(x + \kappa_{t^{'}}((j_1,\hdots,j_{t^{'}})), j, 0)| \mathcal{H}] 
\\
&
= E[G_0(x + \kappa_{t^{'}}((j_1,\hdots,j_{t^{'}})) + \kappa_{t}((j, 0)), (j, 0))| \mathcal{H}]
\\
&
= E[G_0(x + \kappa_{t^{'}}((j_1,\hdots,j_{t^{'}})), (j, 0))| \mathcal{H}] = E[M_0| \mathcal{H}]
\end{split}
\end{align}
almost surely. Since $G$ is continuous the relation established in (\ref{hamv}) above holds for all $x \in [-m,m]$, for all $j \in [-m,m]^{t-1}$, and for almost every $\omega\in\Omega$. The choice $M_1 = E[M_0| \mathcal{H}]$ gives a desired lower bound. 

%We now turn our attention to constructing the bound for $\bar G$. To this end using Assumption \ref{A3} there exists $\bar M = \bar M(m, t^{'})$ so that for $x \in [-m,m]$ we have
%\begin{equation}\label{julikutya}
%\bar G(x) = E[\bar G_0(x + P_{t^{'}}^2 \delta_{t^{'}} / 4)| \mathcal{H}] \geq E[\bar M| \mathcal{H}]
%\end{equation}
%almost surely. Using continuity of $\bar G$ we have that (\ref{julikutya}) holds for all $x \in [-m,m]$ and for almost every $\omega\in\Omega$, and setting $\bar M_1 = E[\bar M| \mathcal{H}]$ finishes the proof.
\end{proof}


%\begin{proof}[Proof of Corollary \ref{corollary::kettes}]
%The statement is a direct consequence of Lemma  \ref{lemma::monotonicity}, and the proof is trivial using continuity of measure , and the $\Omega$-partition formed by the pullback sets $ H^{-1}([m, m+1]), \ m \in \mathbb{Z}$.
%\end{proof}
%\begin{lemma}\label{kettes}
%Let $H$ be an $\mathcal{H}$-measurable random variable that takes values in $I_2$, and assume that the conditions prescribed by Assumption \ref{A2} and Assumption \ref{A3} are met. Then, there exits an $\mathcal{H}$-measurable quantity $K = K(I_1, I_2)$ such that 
%\begin{equation}
%E[V^{\circ}(x, H)| \mathcal{H}] \leq E[V^{\circ}(x, H \indicator{|H|>K})| \mathcal{H}],
%\end{equation}
%for all $x \in I_1$ and for almost every $\omega \in \Omega$.
%\end{lemma}
%
%\begin{proof}
%Observe that, due to Assumption \ref{A2}, for every $\omega \in \Omega$ there exists $K(\omega)$ so that for all $h \in \mathbb{R}$ it holds that
%\begin{align}\label{M}
%\begin{split}
%|h| > K(\omega) \implies \mathbb{E}&[V^{\circ}(l+1, h)| \mathcal{H}](\omega) \leq \mathbb{E}[V^{\circ}(l, 0)| \mathcal{H}](\omega).
%\end{split}
%\end{align}
%Now note, that $\omega \to K(\omega)$ is $\mathcal{H}$-measurable as a random variable. On the event $\{|h| > K\}$, using Lemma \ref{lemma::monotonicity}, and the statement in (\ref{M}), we have 
%\begin{align}\label{boundidstratineq}
%\begin{split}
%&\mathbb{E}[V^{\circ}(x, h)|\mathcal{H}] \leq
%\mathbb{E}[V^{\circ}(l+1, h)|\mathcal{H}],
%\\
%&\leq \mathbb{E}[V^{\circ}(l, 0)|\mathcal{H}] \leq 
%\mathbb{E}[V^{\circ}(x, 0)|\mathcal{H}],
%\end{split}
%\end{align}
%for every $x \in [l, l+1]$, $h \in [m, m+1]$, and almost every $\omega \in \Omega$. That is
%\begin{align}\label{bdStrat}
%\begin{split}
%&\mathbb{E}[V^{\circ}(x, H)|\mathcal{H}] \leq \mathbb{E}[V^{\circ}(x, H \indicator{|H| > K})|\mathcal{H}],
%\end{split}
%\end{align}
%for every $x \in [l, l+1]$, and almost every $\omega \in \Omega$, proving the claim.
%\end{proof}

\section{The generic step, dynamic programming}\label{dyna}

First, in a phase of \emph{bakcward induction} we construct actions that -- depending parametrically on 
previous decisions and accumulated wealth -- are optimal in an instantaneous sense. These actions however would 
only be optimal in one-step markets. 

Then, we use these actions to build a strategy for the entire interval of trading, and this strategy 
will serve as a \emph{candidate strategy} for optimal trading. 

Second, with a \emph{forward iteration} we show that the \emph{candidate} indeed represents a 
series of actions that dominates all admissible executions in terms of expected utility: arriving to the conclusion of the paper. 

Assumptions \ref{uint} and \ref{deltacska} will be in force from now on. Fix $C_{u}\geq 0$ such that $u(x)\leq C_{u}$ for
all $x\in\mathbb{R}$. We denote with $\Xi_{t}$ the $\mathcal{F}_{t}$-measurable random variables.

\begin{proof}[Proof of Theorem \ref{main}]
Define $\tilde{G}_{T}:\Omega\times \mathbb{R}\times \mathbb{R}^{T}\to \mathbb{R}$ as
\begin{equation}\label{okt1}
\tilde{G}_{T}(\omega,x,h_{1},\ldots,h_{T}):=u(x-B(\omega)),\ \omega\in\Omega,\ x,h_{1},\ldots,h_{T}\in\mathbb{R}.
\end{equation}
Note that $h_{1},\ldots,h_{T}$ are dummy variables here, and $\tilde{G}_{T}$ is continuous and nondecreasing in $x$ almost surely. 

The first step of the \emph{backward induction} is different from the other steps
since $h_{T}=-h_{1}-\ldots-h_{T-1}$ due to the constraint on liquidation. 
We thus consider $\tilde{G}_{T-1} : \Omega \times \mathbb{R} \times \mathbb{R}^{T-1} \to \mathbb{R}$ with the definition
\begin{align}\label{okt2}
\begin{split}
& \tilde{G}_{T-1}(\omega,x,h_{1},\ldots,h_{T-1}) \\
&=
E\left[\tilde{G}_{T}\left(x+\kappa_{T}\left(h_{1},\ldots,h_{T-1},-\sum_{k=1}^{T-1}h_{k}\right), h_1, \ldots, h_{T-1},
-\sum_{k=1}^{T-1}h_{k}\right)\left\vert\mathcal{F}_{T-1}\right.\right].
\end{split}
\end{align}

To start the \emph{backward induction} one has to examine whether the conditions prescribed by Assumption \ref{A1}, 
Assumption \ref{A2}, and Assumption \ref{A3} hold with the choice $G_0 = \tilde{G}_{T-1}$. To this end, we note that 
the function $\tilde{G}_{T-1}$ is jointly continuous in its real variables, it is non-decreasing in the first 
real variable almost surely. Using the bound in (\ref{lambda}) we define $\hat G_{T-1}: \Omega \times \mathbb{R} \to \mathbb{R}$ as
$$\hat G_{T-1}(\omega,x) = E\left[u\left(x + \frac{P_T^2(\omega) \delta_T(\omega)}{4}\right)\vert\mathcal{F}_{T-1}\right],$$
and we note that -- due to Assumption \ref{uint} -- for all 
$x \in \mathbb{R}$ it holds that  $\hat G_{T-1}(x) \leq C_{u}$, and as $x \to -\infty$ the quantity $\hat G_{T-1}(x)$ tends to 
$- \infty$ as $x\to -\infty$, in the $\mathbb{P}$-almost sure sense: 
$u$ does so by assumption and due to boundedness from above the reverse Fatou Lemma is applicable. 
Moreover, for all $x, h_1, \ldots, h_{T-1}$, we have $$\tilde{G}_{T-1}(\omega,x,h_{1},\ldots,h_{T-1}) \leq \hat{G}_{T-1}(\omega,x)$$ 
almost surely, leading us to the choice $\bar G_{0} = \hat G_{T-1}$ (again following notation of Section \ref{single}). 

We will establish the following claim after the present proof.

%By Assumption \ref{deltacska}, $$\frac{|\sum_{k=1}^{T-1}h_{k}|}{\delta_{T}}\leq{}
%\frac{|\sum_{k=1}^{T-1}h_{k}|}{\delta},$$ and from Assumption \ref{uint}, for all $x,y\in\mathbb{R}$,
%\begin{equation}
%E[u(x+yP_{T}-B)\vert\mathcal{F}_{T-1}]>-\infty
%\end{equation}
%almost surely. The latter implies also
%\begin{equation}\label{inti1}
%E[u(X+YP_{T}-B)\vert\mathcal{F}_{T-1}]>-\infty.
%\end{equation}
%almost surely, for any $\mathcal{F}_{T-1}$-measurable $X,Y$. Since $\tilde{G}_{T}\leq \tilde{G}_{T-1}$ and \eqref{inti} holds, 
%$\rho_{j,t}\leq 1$ always holds, and the $\kappa_t$ is continuous almost surely, consequently 
\begin{claim}\label{claim} For any $m > 0$  there exists 
an $\mathcal{F}_{T-1}$-measurable and integrable $M = M(m)$ such that for all $1\leq t\leq T-1$ and for all  
$(x,h_1, \ldots, h_{T-1}) \in [-m,m]^{T}$ 
we have $$M(m) \leq G_{T-1}(x + \kappa_{t}(h_1, \ldots, h_{t}), h_1, \ldots, h_{T-1})$$
almost surely.
\end{claim}

One can thus conclude that Assumption \ref{A1}, Assumption \ref{A2}, and Assumption \ref{A3} of Section \ref{single} 
are satisfied and we are ready to perform the first step of the \emph{backward induction}: the 
lemmas of Section \ref{single}, for the first step, will be utilized with the choice $t = T-1$ and $G_0 = \tilde{G}_{T-1}$,
$\bar{G}_{0}:=\hat{G}_{T-1}$.


Lemmas \ref{continuous_version}, \ref{lemma::infinite_loss}, \ref{lemma::monotonicity_}, \ref{lemma::kettes}, 
and \ref{lemma::pointwise_essup} produce a mapping -- denoted by $G$ in the their own context -- 
which, in our notation will take the form $\tilde{G}_{T-2}:\Omega \times \mathbb{R}^{T-1} \to \mathbb{R}$ with the following properties. 
The $\tilde{G}_{T-2}$ is $\mathcal{F}_{T-2} \otimes \mathcal{B}(\mathbb{R}^{T-1})$-measurable, non-decreasing in its first real variable, 
jointly continuous in all its real variables, and for all $x, h_1, \ldots, h_{T-2} \in \mathbb{R}$ it almost surely satisfies
\begin{align}\label{okt3}
\begin{split}
& \tilde{G}_{T-2}(x, h_1, \ldots, h_{T-2})
\\
&=
\mbox{ess.sup}_{H \in \Xi_{T-2}}E\left[\tilde{G}_{T-1}(x + \kappa_{T-1}(h_1, \ldots, h_{T-2}, H), 
h_1, \ldots, h_{T-2}, H) \vert \mathcal{F}_{T-2} \right].
\end{split}
\end{align}
Furthermore, Lemmas \ref{lemma::harmas}, and \ref{optimal_H} ensure the existence of a mapping 
$H^{*}_{T-1}:\Omega \times \mathbb{R}^{T-1} \to \mathbb{R}$ that is $\Omega \otimes \mathcal{B}(\mathbb{R}^{T-1})$-measurable, 
and is such that for all random variables $X, H_1, \ldots, H_{T-2}$ that are measurable with respect to  
$\mathcal{F}_{T-2}$, with the notation $H^{*}_{T-1}= H^{*}_{T-1}(X, H_1, \ldots, H_{T-2})$, we have
\begin{align}\label{okt4}
\begin{split}
& \tilde{G}_{T-2}(X, H_1, \ldots, H_{T-2})
\\
&=
E\left[\tilde{G}_{T-1}(X + \kappa_{T-1}(H_1, \ldots, H_{T-2}, H^{*}_{T-1}), H_1, \ldots, H_{T-2}, H^{*}_{T-1}) \vert 
\mathcal{F}_{T-1} \right],
\end{split}
\end{align}
$\mathbb{P}$-almost surely.


Lemmas \ref{x_bound}, and \ref{lemma::lower_bd} imply that there exists an action-independent bound 
$\hat G_{T-2}$, with properties analogous to the $\hat G_{T-1}$ presented above, and thus finally, 
we arrive to the conclusion that the quantity $\tilde{G}_{T-2}$ is such that it again satisfies 
Assumption \ref{A1}, Assumption \ref{A2}, and Assumption \ref{A3}. That is, in 
the next step of iteration, the choice $G_0 = \tilde{G}_{T-2}$, $\bar{G}_{0}:=\hat{G}_{T-2}$ can be made.

Iterating backwards in this manner goes with the usual mechanics of induction. Take $\tilde{G}_{T-2}$ 
as a starting point. 

When $\tilde{G}_{t^{'}}, \ldots, \tilde{G}_{T}$ 
(and $H^*_{t^{'}-1}, \ldots, H^*_{T-1}$) are given for some $t^{'} \leq T-2$, applying the lemmas of Section  
\ref{single} with the choice $G_{0} = \tilde{G}_{t^{'}}$, $\bar{G}_{0}:=\hat{G}_{t'}$ and $t = t^{'}$ 
yield $\tilde{G}_{t^{'}-1}$, and with this procedure we construct the pairs
\begin{equation}\label{pairs}
(H^{*}_{T-1}, \tilde{G}_{T-2}), (H^{*}_{T-2}, \tilde{G}_{T-3}), \ldots, (H^{*}_1,\tilde{G}_{0})
\end{equation}
with the properties shown below.


For $t \in \{0,1, \ldots, T-2\}$, $\tilde{G}_{t}:\Omega \times \mathbb{R}^{t+1} \to \mathbb{R}$ is 
$\mathcal{F}_{t} \otimes \mathcal{B}(\mathbb{R}^{t+1})$-measurable, non-decreasing in its first 
real variable, jointly continuous in its real variables, and for all $x, h_1, \ldots, h_{t} \in \mathbb{R}$, 
in the almost sure sense we have
\begin{align*}
\begin{split}
\tilde{G}_{t}(x, h_1,& \ldots, h_{t})
\\
=&
\mbox{ess.sup}_{H \in \Xi_{t}}E\left[\tilde{G}_{t+1}(x + \kappa_{t+1}(h_1,\ldots,h_{t},H), h_1, \ldots, h_{t}, H) \vert \mathcal{F}_t \right].
\end{split}
\end{align*}
The mapping $H^{*}_{t+1}:\Omega \times \mathbb{R}^{t} \to \mathbb{R}$ is $\Omega \otimes \mathcal{B}(\mathbb{R}^{t})$-measurable 
and for all $\mathcal{F}_{t}$-measurable random variables $X, H_1, \ldots, H_{t}$ we have for 
$\tilde H^{*}_{t+1}= H^{*}_{t+1}(X, H_1, \ldots, H_{t})$ that
\begin{align}\label{zenta}
\begin{split}
\tilde{G}_{t}&(X, H_1, \ldots, H_{t}) 
\\
=&
E\left[\tilde{G}_{t+1}(X + \kappa_{t+1}(H_1,\ldots,H_t, \tilde H^*_{t+1}), H_1, \ldots, H_{t}, 
\tilde H^{*}_{t+1}) \vert \mathcal{F}_t \right],
\end{split}
\end{align}
holds in the $\mathbb{P}$-almost sure sense. 


Introducing notation, for any admissible trading strategy $H = (H_t)_{t \in \{1, \ldots, T \} }$, we denote by 
$\Gamma_t( H)$ the strategy $( H_s)_{s \in \{1, \ldots, t\}}$, the same trading strategy as $H$, but without liquidation, 
and corresponding to the trading interval up to $t$.

Now we construct the \emph{candidate strategy} using the mapping in (\ref{pairs}). Let $\hat H_1 = H^{*}_1(z)$. 
We define the optimal steps using a forward recursion. That is, when the $\hat H_t$ is constructed for 
some $T -1 \geq t \geq 1$, with accumulated wealth $\hat X_{t} = \xi_{t}^{\Gamma_t(\hat H)}$, the 
next action $\hat H_{t+1}$ is given as
\begin{equation}\label{anyulontul}
\hat H_{t+1} = H^{*}_{t+1}(\hat X_t, \hat H_{1}, \ldots, \hat H_{t}).
\end{equation}
Then, we set $\hat H_T = -\hat H_{T-1} - \ldots - \hat H_1$.


We proceed with the \emph{forward iteration}, giving an upper bound for all expected payoffs which -- 
as we shell see -- equals the expected payoff associated with the \emph{candidate strategy} $\hat H$.


%let us introduce a companion strategy defined as  $\mathcal{R}_{\tau}(H) = \{ H_1, \ldots H_{\tau-1}, \hat H_{\tau } , \ldots , \hat H_{T-1}, \hat H_{T}^{'}  \}$, where $\hat H_{T}^{'} = - H_1 - \cdots - H_{\hat T -1} - \hat H_{\hat T} - \cdots - \hat H_{T-1}$.
Letting $H_1, \ldots, H_{T-1}, H_T$ denote an arbitrary series of admissible actions. Iterating with the rule in (\ref{zenta}) yields
\begin{align}\label{arvaszamar}
\begin{split}
&E[u(z + \xi_T^{H} + B)]
\\
&=
E[ E[u(z + \xi_{T-1}^{\Gamma_{T-1}(H)} + \kappa_T(H_1, \ldots, H_{T-1}, - \sum_{k = 1}^{T-1} H_{k}) + B)] \vert \mathcal{F}_{T-1}]]
\\
&=
E[ \tilde{G}_{T-1}(z + \xi_{T-2}^{\Gamma_{T-2}(H)} + \kappa_{T-1}(H_1, \ldots, H_{T-1}), H_1, \ldots, H_{T-1}) ]
\\
& 
\leq
E[ E[ \tilde{G}_{T-2}(z + \xi_{T-3}^{\Gamma_{T-3}(H)} + \kappa_{T-2}(H_1, \ldots, H_{T-2}), H_1, \ldots, H_{T-2}) \vert \mathcal{F}_{T-2} ] ]
\\
& \vdots
\\
&\phantom{} \leq  E[\tilde{G}_0(z)].
\end{split}
\end{align}


Furthermore, as a result of (\ref{okt1}), (\ref{okt2}), (\ref{okt3}), and (\ref{okt4}) we have that 
\begin{align}\label{csapzottmacska}
\begin{split}
\tilde{G}_0(z) =& E[ \tilde{G}_1(z + \kappa_1(\hat H_1), \hat H_1)]
\\
=&  E[\tilde{G}_2(z + \kappa_1(\hat H_1) + \kappa_2(\hat H_1, \hat H_2), \hat H_1, \hat H_2)]
\\
\vdots&
\\
=& E[\tilde{G}_{T-1}(z + \sum_{i = 1}^{T-1}\kappa_{i}(\hat H_1, \ldots, \hat H_{i}), \hat H_1, \ldots, \hat H_{T-1})]
\\
=& E[u(z + \xi_T^{\hat H} + B)].
\end{split}
\end{align}
Thus, we can conclude, due to (\ref{csapzottmacska}) and (\ref{arvaszamar}), that the candidate 
strategy given with trading actions $(\hat H_t)_{t = 1, \ldots, T}$, defined in (\ref{anyulontul}), is indeed optimal. 
Theorem \ref{main} is now shown.
\end{proof}

\begin{proof}[Proof of Claim \ref{claim}]
For $h_{1},\ldots,h_{T-1}\in [-m,m]$ we clearly have
$$
\kappa_{t}(h_{1},\ldots,h_{t})\geq -m\zeta_{0}-tm^{2}/\delta+\min\{mP_{t},-mP_{t}\}.	
$$
Similarly,
\begin{eqnarray*} & & 
\kappa_{T}\left(h_{1},\ldots,h_{T-1},-\sum_{k=1}^{T-1}h_{k}\right)\\
&\geq& -(T-1)m\zeta_{0}-(T-1)m^{2}(2T-2)/\delta+\min\{m(T-1)P_{T},-m(T-1)P_{T}\}
\end{eqnarray*}
Define the constant
$$
D_{m}:=-(2T-2)\zeta_{0}-(T-1)m^{2}(2T-1)/\delta.
$$
We can thus see that
$$
\tilde{G}_{T-1}(x+\kappa_{t}(h_{1},\ldots,h_{t}),h_{1},\ldots,h_{T-1})\geq \mathbb{E}\left[{}
\min\{J_{1},J_{2},J_{3},J_{4}\}\left\vert\mathcal{F}_{T-1}\right.\right]
$$
where
\begin{eqnarray*}
J_{1} &=& u(x+D_{m}+mP_{t}+m(T-1)P_{T}-B),\\
J_{2} &=& u(x+D_{m}-mP_{t}+m(T-1)P_{T}-B),\\
J_{3} &=& u(x+D_{m}+mP_{t}-m(T-1)P_{T}-B),\\
J_{4} &=& u(x+D_{m}-mP_{t}-m(T-1)P_{T}-B).	
\end{eqnarray*}
Now \eqref{inti} implies our statement.
\end{proof}



%
%
%%We write out the conditional expectation explicitly:
%%$$
%%E\left[G_{T}\left(x+P_{T}\left(\sum_{k=1}^{T-1}h_{k}\right)
%%-|\sum_{k=1}^{T-1}h_{k}|\left(\rho_{0,T}\zeta_{0}+ \frac{|\sum_{k=1}^{T-1}h_{k}|}{\delta_{T}}+
%%\sum_{j=1}^{T-1} \frac{\rho_{j,T}|h_{j}|}{\delta_{j}}\right)\right) \vert\mathcal{F}_{T-1}\right].
%%$$
%%Here we have used that $\rho_{T,T}=1$. Note that everything is $\mathcal{F}_{t-1}$-measurable here, except
%%$P_{T}$ and $\delta_{T}$.
%
%The function 
%$$
%(x,h_{1},\ldots,h_{T-1})\to G_{T}\left(x+P_{T}\left(\sum_{k=1}^{T-1}h_{k}\right)
%-|\sum_{k=1}^{T-1}h_{k}|\left(\rho_{0,T}\zeta_{0}+ \frac{|\sum_{k=1}^{T-1}h_{k}|}{\delta_{T}}
%+\sum_{j=1}^{T-1} \frac{\rho_{j,T}|h_{j}|}{\delta_{j}}\right)\right)
%$$
%is almost surely continuous, it is bounded from above, so \eqref{inti} and the usual lemma on the construction of
%continuous versions guarantees that there are versions of the conditional expectation such that, outside a $0$-set:
%$$
%G_{T-1}(\cdot,x,h_{1},\ldots,h_{T-1})
%$$
%is a continuous function. What's more, it is non-decreasing in $x$.
%
%From this step on, we apply a generic argument using the one-step case for $G_{t+1}$ and getting optimal strategies
%and regular version for $G_{t}$. Finally, we construct $h_{1}$ and $G_{0}$.
%
%
%%---------------------------------------------------------------
%%---------------------------------------------------------------------
%%--------------------------------------------------------------------------------
%%-----------------S T A T E M E N T S O F T H E M A I N L E M M A A B O V E ------- %-------------------------------------------------------------------------------
%%--------------------------------------------------------------------------
%%------------------------------------------------------------------
%%---------------------------------------------------------------
%
%
%%\begin{lemma}\label{project} For $1\leq t\leq T$,
%%there exist $\mathcal{F}_{t-1}\otimes\mathcal{B}(\mathbb{R})\otimes\mathcal{B}(\mathbb{R})^{t-1}$-measurable
%%functions $G_{t-1}:\Omega\times\mathbb{R}\times\mathbb{R}^{t-1}\to \mathbb{R}$ such that\footnote{We mean that
%%$G_{0}:\Omega\times\mathbb{R}\to\mathbb{R}$.} for
%%almost all $\omega\in\Omega$, $G_{t-1}(\omega,\cdot,\cdot)$ is continuous, it is non-decreasing in its
%%first variable and 
%%\begin{align*}
%%G_{t-1}(x,h_{1},\ldots,h_{t-1})& \\ = \ \ \ \mathrm{ess.}\sup_{H\in \Xi_{t-1}}&
%%\mathbb{E}_{t-1}[G_{t}(x+\kappa_{t}(h_{1},\ldots,h_{t-1},H),h_{1},\ldots,h_{t-1},H)]
%%\end{align*}
%%almost surely, for all $x,h_{1},\ldots,h_{t-1}$. 
%%For all $h_{1},\ldots,h_{t-1}\in\mathbb{R}$,
%%\begin{equation}\label{-infty}
%%\lim_{x\to -\infty}G_{t-1}(x,h_{1},\ldots,h_{t-1})\to -\infty{}
%%\end{equation} 
%%almost surely.
%%There exist integrable random variables $M_{t-1}(n)$, $n\in\mathbb{Z}$ such that
%%\begin{equation}\label{gm}
%%G_{t-1}(n)\geq M_{t-1}(n)	
%%\end{equation} 
%%almost surely.
%%\end{lemma}
%%
%%\begin{proof}
%%We follow an argument of induction. Assume that for a fixed $t \leq T$ we have that $G_t$ is continuous, non-decreasing in its first variable, and for all $x$ and $h_1,...,h_{t-1}$ let $F(x, h_1, ..., h_{t-1})$ denote a version of the essential supremum, that is let
%%\begin{align*}
%%F(x, h_1, ..., h_{t-1}) & \\
%% = \ \ \ \mathrm{ess.}\sup_{H\in \Xi_{t-1}}&
%%\mathbb{E}_{t-1}[G_{t}(x+\kappa_{t}(h_{1},\ldots,h_{t-1},H),h_{1},\ldots,h_{t-1},H)].
%%\end{align*}
%%Also, define $$G_{t-1}(x,h_1,...,h_{t-1}) =\inf_{y>x, y \in \mathbb{Q}} F(y, h_1, ...,h_{t-1}).$$ Outside a negligible set $N \subset \Omega$ the function $F$ restricted to $\mathbb{Q}$ is non-decreasing. Thus, $G_{t-1}$ - outside the negligible set $N$ - is right continuous and non-decreasing by definition. In the followings we will show that it is a version of the essential supremum in consideration. Fix $x \in \mathbb{R}$. On one hand, for $x < q$ we have $$F(x, h_1,...,h_{t-1}) < F(q, h_1,...,h_{t-1})$$ a.s., and in the limit this shows that $F(x, h_1,...,h_{t-1}) \leq G_{t-1}(x, h_1,...,h_{t-1})$. On the other hand, first let $1 \geq q_k, \ k\in \mathbb{N}$, be a sequence in $\mathbb{Q}$ tending to zero, and for each $k$ select $H_k$ so that 
%%\begin{align*}
%% F(x + q_k, h_1,...,&h_{t-1}) - 1/k \\ 
%% \leq \ \ \ & E_{t-1}\left[ G_t(x + q_k + \kappa_t(h_1,...,h_{t-1}, H_k), h_1,...,h_{t-1},H_k) \right]
%%\end{align*}
%%
%%\end{proof}
%
%
%\begin{lemma}\label{opti} For every $H_{j}\in \Xi_{j-1}$, $j=1,\ldots,t-1$ and for all $V_{t-1}\in \Xi_{t-1}$
%there exists $H^{\dagger}_{t}\in \Xi_{t-1}$ such that
%$$
%G_{t-1}(V_{t-1},H_{t-1},\ldots,H_{1})=\mathbb{E}_{t-1}
%[G_{t}(V_{t-1}+\kappa_{t}(H^{\dagger}_{t},H_{t-1},\ldots,H_{1}),H_{t-1},\ldots,H_{1})].
%$$
%\end{lemma}
%
%
%
%
%
%	
%\begin{lemma}
%Let $\hat T \in \{ 1, \hdots, T-1 \}$. For every $V_{T-2}$ for every $T-2$-tuple of random variables $H_1,\hdots, H_{T-2}$ there exists 
%\end{lemma}
%
%
%\begin{proof}[Proof of Theorem \ref{main}.]	
%We apply Lemma \ref{opti} first with the choice $t:=1$,
%$V_{0}:=x$. Denote by $H_{1}^{*}$ the resulting $H_{1}^{\dagger}$.
%Continuing recursively, we apply Lemma \ref{opti} with the choice $V_{t-1}:=\xi_{t-1}^{H^{*}}$,
%$H_{1}:=H_{1}^{*},\ldots,H_{t-1}:=H_{t-1}^{*}$ and denote the resulting $H_{t}^{\dagger}$ by $H_{t}^{*}$.
%We construct a stragegy $H_{t}^{*}$, $1\leq t\leq T$ in this way. We claim that $H^{*}$ is optimal
%for the utility maximization problem.
%
%Indeed, for an arbitrary $H\in\mathcal{A}_{0}$, we have\ldots{}
%As there are equalities everywhere for $H^{*}=H$, we conclude that $H^{*}$ is indeed what we have been looking for.
%\end{proof}
%
%\begin{corollary}\label{ug}
%For all $z\in\mathbb{R}$, $G_{0}(z)=\bar{u}(z)$.	
%\end{corollary}

%\section{Auxiliary Lemmas}\label{aux}


\section{Conclusion}\label{conclu}

One could add a solvency constraint (that is, $z+\xi^{X}_{t}\geq 0$ for all $t$) with minimal modifications
in the arguments. Setting $B:=0$, utility maximization for $u$ defined on the positive real axis 
could be treated in this way.{}

It is unclear whether (and how) results of the present paper could be transferred to continuous-time models.
We rely, in an essential way, on the fact that the treated portfolio optimization problem can be broken down into one-step
problems. In continuous time such an approach is out of question.



\medskip{}

{\footnotesize \noindent\textbf{Acknowledgments.} Both authors gratefully 
acknowledge the support of 
the National Research, Development and Innovation Office (NKFIH) through grant K 143529 
and also within the framework of the Thematic Excellence Program 2021 (National Research subprogramme 
``Artificial intelligence, large networks, data security: mathematical foundation and applications'').
The second author also thanks for the support of NKFIH grant KKP 137490.}


%\bibliographystyle{plain}
%\bibliography{transient_impact}


\begin{thebibliography}{9}

\bibitem{peter-yan} P. Bank and Y. Dolinsky.
\newblock Continuous-time duality for super-replication with transient price impact.
\newblock \emph{Ann. Appl. Probab.}, 29:3893--3917, 2019.  

\bibitem{pym} P. Bank, Y. Dolinsky, M. R\'asonyi.{}
\newblock What if we knew what the future brings? Optimal investment for a frontrunner with price impact.
\newblock \emph{Appl. Math. Opt.}, vol. 86, paper no. 25, 1--24, 2022.

\bibitem{bank-voss} P. Bank and M. Vo\ss{}.
\newblock Optimal investment with transient price impact.
\newblock \emph{SIAM J. Finan. Math.}, 10:723--768, 2019. 

\bibitem{laurence} L. Carassus and M. R\'asonyi.{}
\newblock Maximization of non-concave utility functions in
discrete-time financial market models. 
\newblock \emph{Math. Oper. Res.}, 41:146--173, 2016.

%%CASTAING , C. and VALADIER , M. (1977). Convex Analysis and Measurable Multifunctions.
%%Lecture Notes in Math. 580. Springer, Berlin

%\bibitem{dm}
%C.~Dellacherie and P.-A. Meyer.
%\newblock \emph{Probability and potential}.
%\newblock North-Holland, Amsterdam, 1979.

%\bibitem{nev}
%J. F. C. Kingman.
%\newblock Discrete-Parameter Martingales.
%\newblock \emph{Royal Statistical Society}. Series A: General, Volume 139, Issue 4, July 1976, Pages 547–548, https://doi.org/10.2307/2344365


\bibitem{gr} P. Guasoni and M. R\'asonyi.{}
\newblock Hedging, arbitrage and optimality under superlinear friction.
\newblock \emph{Ann. Appl. Probab.}, 25:2066--2095, 2015.

\bibitem{ks} Yu. M. Kabanov and Ch. Stricker. 
\noindent A teacher’s note on no-arbitrage criteria.
\noindent \emph{In:
Séminaire de Probabilités, vol. XXXV}, 149--152, Springer-Verlag, 2001.

%\bibitem{bounded} M. R\'asonyi. 
%\newblock Optimal investment with nonconcave utilities in discrete-time markets.
%\emph{SIAM J. Finan. Math.}, 6:517--529, 2015.

\bibitem{neveu}
{J. Neveu} 
\newblock {\em Discrete parameter martingales.}
\newblock North-Holland, Amsterdam,1975.

\bibitem{lukasz} M. R\'asonyi and L. Stettner.
\newblock On utility maximization
in discrete-time market models.
\newblock {\em Ann. Appl. Probab.}, 15:1367--1395, 2005.

\end{thebibliography}

\end{document}

Reviewer: Faustino Gozzi, Klaus Reiner Schenk-Hoppé, Mihai Sirbu, Xiaolu Tan, Nizar Touzi
